%% 48장 Gaussian Elimination and LU Factorization — Theory
\chapter{Gaussian Elimination and LU Factorization}
\label{ch:lu}

\abstract{가우스 소거는 행렬을 단위 하삼각 $L$과 상삼각 $U$의 곱으로 나누는 일이다. 이 장은 LU 분해의 존재와 부분 피벗의 안정성 이론에서 시작해 블록 알고리즘과 재귀 알고리즘으로 가고, stardust의 \code{lu}, \code{solve}, \code{lu\_solve}, \code{inv}, \code{det}, \code{slogdet}이 LAPACK을 감싸 그것을 어떻게 구현하는지까지 따라간다.}

\section{Elimination as Factorization}
\label{sec:lu-factorization}

\subsection{One Step of Elimination}

$n\times n$ 행렬 $A$로 $Ax=b$를 푼다고 하자. 가우스 소거의 첫 단계는 첫 행의 배수를 아래 행들에서 빼서 첫 열의 대각 아래를 $0$으로 만든다. $a_{11}\ne0$이면 \term{승수}{multiplier} $\ell_{i1}=a_{i1}/a_{11}$($i=2,\dots,n$)를 쓰고 이 연산은 행렬
\[
M_1=I-\ell_1e_1^{\mathsf T},\qquad \ell_1=(0,\ell_{21},\dots,\ell_{n1})^{\mathsf T}
\]
을 왼쪽에서 곱하는 것과 같다. $A$를 첫 행과 첫 열로 나누어 쓰면 이 단계가 무엇을 남기는지가 분명해진다.
\begin{equation}
\label{eq:lu-schur-step}
A=\begin{pmatrix}a_{11}&a_{12}^{\mathsf T}\\a_{21}&A_{22}\end{pmatrix}
=\begin{pmatrix}1&0\\a_{21}/a_{11}&I\end{pmatrix}
\begin{pmatrix}a_{11}&a_{12}^{\mathsf T}\\0&S\end{pmatrix},\qquad
S=A_{22}-\frac{a_{21}a_{12}^{\mathsf T}}{a_{11}}.
\end{equation}
오른쪽의 곱을 직접 전개하면 왼쪽이 된다. $(2,2)$ 블록이 $(a_{21}/a_{11})a_{12}^{\mathsf T}+S=A_{22}$이기 때문이다. $S$를 $a_{11}$에 대한 \term{Schur 보수}{Schur complement}라 한다. 소거의 한 단계는 승수 열 하나를 $L$에 적고, 첫 행을 $U$에 적고, 남은 $(n-1)\times(n-1)$ 문제를 $S$로 바꾸는 일이다(\cref{fig:lu-step}). $S$의 원소 하나를 고치는 데 곱셈 하나와 뺄셈 하나가 드므로 한 단계의 비용은 $2(n-1)^2$ 플롭에 나눗셈 $n-1$번이다.

\begin{figure}[t]
\centering
\begin{tikzpicture}[x=0.36cm,y=0.36cm,font=\small]
  % before
  \draw (0,0) rectangle (6,6);
  \fill[gray!45] (0,5) rectangle (1,6);
  \fill[gray!20] (1,5) rectangle (6,6);
  \fill[gray!20] (0,0) rectangle (1,5);
  \node at (0.5,5.5) {$a_{11}$};
  \node at (3.5,5.5) {$a_{12}^{\mathsf T}$};
  \node at (0.5,2.5) {$a_{21}$};
  \node at (3.5,2.5) {$A_{22}$};
  \node at (3,-0.9) {$A$};
  \draw[->,thick] (6.8,3) -- (8.6,3);
  % after: L column, U row, S
  \begin{scope}[xshift=8.1cm]
  \draw (0,0) rectangle (6,6);
  \fill[gray!45] (0,5) rectangle (1,6);
  \fill[gray!45] (1,5) rectangle (6,6);
  \fill[pattern=north east lines] (0,0) rectangle (1,5);
  \draw (0,0) rectangle (1,5);
  \fill[gray!10] (1,0) rectangle (6,5);
  \node at (0.5,5.5) {$u_{11}$};
  \node at (3.5,5.5) {$u_{12}^{\mathsf T}$};
  \node[fill=white,inner sep=1pt] at (0.5,2.5) {$\ell_{1}$};
  \node at (3.5,2.5) {$S$};
  \node at (3,-0.9) {한 단계 뒤 (같은 자리)};
  \end{scope}
\end{tikzpicture}
\caption{소거의 한 단계. 첫 행은 그대로 $U$의 첫 행이 되고, 첫 열의 대각 아래는 승수 $\ell_1=a_{21}/a_{11}$로 바뀌며, 나머지는 Schur 보수 $S=A_{22}-\ell_1u_{12}^{\mathsf T}$로 갱신된다. 구현은 이 셋을 모두 $A$의 자리에 덮어쓴다.}
\label{fig:lu-step}
\end{figure}

$S$에 같은 일을 되풀이하면 $n-1$번의 단계 뒤에 상삼각 행렬 $U=M_{n-1}\cdots M_1A$가 남는다. 이것을 $A$에 대해 풀면 $A=M_1^{-1}\cdots M_{n-1}^{-1}U$이고, 다음 보조정리는 앞의 곱이 승수를 제자리에 늘어놓은 단위 하삼각 행렬이라는 것을 보인다.

\begin{lemma}
\label{lem:lu-multipliers}
$\ell_k$의 처음 $k$개 성분이 $0$일 때 $M_k=I-\ell_ke_k^{\mathsf T}$의 역은 $I+\ell_ke_k^{\mathsf T}$이고
\[
M_1^{-1}M_2^{-1}\cdots M_{n-1}^{-1}=I+\sum_{k=1}^{n-1}\ell_ke_k^{\mathsf T}
\]
이다. 곧 이 곱의 $(i,k)$ 원소($i>k$)는 $k$번째 단계의 승수 $\ell_{ik}$이다.
\end{lemma}

\begin{proof}
$e_k^{\mathsf T}\ell_k=0$이므로 $(I-\ell_ke_k^{\mathsf T})(I+\ell_ke_k^{\mathsf T})=I-\ell_k(e_k^{\mathsf T}\ell_k)e_k^{\mathsf T}=I$이다. 곱에 대해서는 $j<k$이면 $e_j^{\mathsf T}\ell_k=0$이라는 것을 쓴다. $(I+\ell_1e_1^{\mathsf T})\cdots(I+\ell_{n-1}e_{n-1}^{\mathsf T})$를 전개하면 둘 이상의 $\ell_je_j^{\mathsf T}$가 곱해진 항은 모두 $\ell_{j}e_{j}^{\mathsf T}\ell_{k}e_k^{\mathsf T}$($j<k$) 꼴의 인수를 품는다. 그런데 $e_j^{\mathsf T}\ell_k=0$이므로 그런 항은 모두 $0$이고 일차 항만 남는다.
\end{proof}

순서가 중요하다는 점을 짚어 둔다. 역순의 곱 $M_{n-1}\cdots M_1$은 승수가 깔끔하게 늘어서지 않는다. $L$을 따로 계산할 필요 없이 승수를 $A$의 대각 아래에 적어 두기만 하면 된다는 것이 위 보조정리의 실질적인 뜻이고, 모든 LU 구현이 $L$과 $U$를 $A$ 하나의 자리에 덮어쓰는 까닭이다.

\subsection{Existence and Uniqueness}

소거가 끝까지 가려면 단계마다 대각 원소, 곧 \term{피벗}{pivot}이 $0$이 아니어야 한다. 그 조건을 $A$로 말하면 다음과 같다. $A_k$는 $A$의 왼쪽 위 $k\times k$ 부분행렬, 곧 \term{선행 주부분행렬}{leading principal submatrix}이다.

\begin{theorem}
\label{thm:lu-existence}
$A\in\mathbb C^{n\times n}$가 단위 하삼각 $L$과 정칙 상삼각 $U$의 곱 $A=LU$로 쓰일 필요충분조건은 모든 $k=1,\dots,n$에서 $A_k$가 정칙인 것이다. 이때 분해는 하나뿐이고 $u_{kk}=\det A_k/\det A_{k-1}$($\det A_0=1$)이다.
\end{theorem}

\begin{proof}
$A=LU$이면 블록 곱으로 $A_k=L_kU_k$이다. $L$과 $U$가 삼각이어서 왼쪽 위 블록의 곱에 나머지 블록이 끼어들지 않기 때문이다. 그러므로 $\det A_k=\det L_k\det U_k=u_{11}\cdots u_{kk}\ne0$이고, 이웃한 두 $k$의 몫이 $u_{kk}=\det A_k/\det A_{k-1}$이다.

역은 $n$에 대한 귀납법이다. $n=1$은 $L=1$, $U=a_{11}$이다. $n>1$이면 $A_{n-1}$의 모든 선행 주부분행렬이 정칙이므로 귀납 가정으로 $A_{n-1}=L_{n-1}U_{n-1}$이다. 이제
\[
A=\begin{pmatrix}A_{n-1}&c\\r^{\mathsf T}&\alpha\end{pmatrix}
=\begin{pmatrix}L_{n-1}&0\\ \ell^{\mathsf T}&1\end{pmatrix}
\begin{pmatrix}U_{n-1}&u\\0&\beta\end{pmatrix}
\]
가 되도록 $L_{n-1}u=c$, $\ell^{\mathsf T}U_{n-1}=r^{\mathsf T}$, $\beta=\alpha-\ell^{\mathsf T}u$로 놓는다. $L_{n-1}$과 $U_{n-1}$이 정칙이므로 $u$와 $\ell$은 하나로 정해지고, $\det A=\det U_{n-1}\,\beta\ne0$에서 $\beta\ne0$이다.

유일성: $L_1U_1=L_2U_2$이면 $L_2^{-1}L_1=U_2U_1^{-1}$이다. 왼쪽은 단위 하삼각이고 오른쪽은 상삼각이므로 둘 다 대각이어야 하며, 왼쪽의 대각이 모두 $1$이므로 단위행렬이다.
\end{proof}

이 정리는 피벗 없이 소거할 수 있는 행렬을 정확히 가려내지만 계산에는 쓸모가 적다. $A_k$가 정칙이어도 거의 특이하면 피벗이 작아지고, 다음 절에서 보듯 작은 피벗은 반올림 오차를 크게 키운다. 정칙이 아닌 $A$나 $m\ne n$인 직사각 행렬에도 소거는 그대로 할 수 있다. $m\times n$ 행렬을 $p=\min(m,n)$번 소거하면 $m\times p$ 단위 하사다리꼴 $L$과 $p\times n$ 상사다리꼴 $U$가 나오고, 피벗이 $0$인 단계에서는 그 열을 건너뛰면($\ell_k=0$) 분해는 여전히 존재한다. 이때 $U$의 대각에 $0$이 나타난다. stardust의 \code{lu}가 특이 행렬에도 에러를 내지 않는 것이 이 사실에 기대고 있다(\cref{sec:lu-stardust-family}).

\section{Pivoting and Growth}
\label{sec:lu-pivoting}

\subsection{Why Pivoting Is Needed}

피벗이 정확히 $0$이 아니어도 작으면 문제가 생긴다. $\epsilon=10^{-20}$일 때
\[
A=\begin{pmatrix}\epsilon&1\\1&1\end{pmatrix},\qquad b=\begin{pmatrix}1\\2\end{pmatrix}
\]
의 해는 $x_1=1/(1-\epsilon)$, $x_2=(1-2\epsilon)/(1-\epsilon)$로 둘 다 $1$에 매우 가깝다. 피벗 없이 소거하면 승수는 $\ell_{21}=10^{20}$이고 $u_{22}=1-10^{20}$은 binary64에서 $-10^{20}$으로 반올림된다. 뒤로 대입하면 $\hat x_2=\mathrm{fl}((2-10^{20})/(1-10^{20}))=1$은 맞지만 $\hat x_1=\mathrm{fl}((1-\hat x_2)/\epsilon)=0$이다. 오차는 $100\%$이다. 잃어버린 것은 $u_{22}$를 만들 때 $a_{22}=1$이 $10^{20}$에 묻혀 사라진 정보이다. 계산된 $\hat L\hat U$는 $\begin{pmatrix}\epsilon&1\\1&0\end{pmatrix}$이어서 $A$와 $(2,2)$ 원소가 $1$만큼 다르다. 이 차이는 $A$의 원소 크기만큼 크다. 반면 두 행을 바꾸고 소거하면 승수가 $\epsilon$이고 $u_{22}=1-\epsilon$이 정확하게 남아 $\hat x$가 해와 반올림 수준에서 같다.

\subsection{Partial Pivoting}

\term{부분 피벗}{partial pivoting}은 $k$번째 단계에서 $k$열의 대각과 그 아래 원소 가운데 절댓값이 가장 큰 것을 찾아 그 행을 $k$번째 행과 바꾼 뒤에 소거한다. 그러면 모든 승수가 $|\ell_{ik}|\le1$이다. 행 교환을 순열 행렬 $P_k$로 쓰면 소거는 $U=M_{n-1}P_{n-1}\cdots M_1P_1A$이다.

\begin{theorem}
\label{thm:lu-partial-pivoting}
모든 $A\in\mathbb C^{n\times n}$에 대해 부분 피벗 소거는 $PA=LU$를 준다. $P$는 순열 행렬, $L$은 $|\ell_{ij}|\le1$인 단위 하삼각, $U$는 상삼각이고, $A$가 정칙이면 $U$도 정칙이다. $L$의 $k$열은 $k$번째 단계의 승수를 그 뒤의 교환 $P_{n-1}\cdots P_{k+1}$로 섞은 것이다.
\end{theorem}

\begin{proof}
$k$번째 단계에서 $k$열의 대각 아래가 모두 $0$이면 교환 없이 $\ell_k=0$으로 둔다($M_k=I$). 그렇지 않으면 가장 큰 원소를 고르므로 $|\ell_{ik}|\le1$이다. 이제
\[
\Pi_k=P_{n-1}\cdots P_{k+1}\quad(\Pi_{n-1}=I),\qquad\tilde M_k=\Pi_kM_k\Pi_k^{-1}
\]
로 두자. $P_j$($j>k$)는 행 $j$와 그 아래의 한 행을 바꾸므로 $\Pi_k$는 $k$ 이하의 행을 건드리지 않고 $\Pi_ke_k=e_k$, $e_k^{\mathsf T}\Pi_k^{-1}=(\Pi_ke_k)^{\mathsf T}=e_k^{\mathsf T}$이다(순열 행렬은 직교한다). 그러므로 $\tilde M_k=I-(\Pi_k\ell_k)e_k^{\mathsf T}$이고 $\tilde\ell_k=\Pi_k\ell_k$는 처음 $k$개 성분이 $0$인 채로 성분만 섞인 벡터이다. 이제 $\Pi_{k-1}=\Pi_kP_k$이므로 $\Pi_k^{-1}\Pi_{k-1}=P_k$이고
\begin{align*}
\tilde M_{n-1}\cdots\tilde M_1\,\Pi_0
&=\Pi_{n-1}M_{n-1}(\Pi_{n-1}^{-1}\Pi_{n-2})M_{n-2}\cdots M_1(\Pi_1^{-1}\Pi_0)\\
&=M_{n-1}P_{n-1}\cdots M_1P_1
\end{align*}
이다. 따라서 $P=\Pi_0=P_{n-1}\cdots P_1$, $L=\tilde M_1^{-1}\cdots\tilde M_{n-1}^{-1}$로 두면 $\tilde M_{n-1}\cdots\tilde M_1PA=U$, 곧 $PA=LU$이고 \cref{lem:lu-multipliers}로 $L$은 섞인 승수 $\tilde\ell_k$를 열로 가진 단위 하삼각이다. 섞어도 크기는 그대로이므로 $|\ell_{ij}|\le1$이다. $A$가 정칙이면 각 단계에서 $k$열의 대각과 그 아래가 모두 $0$일 수 없으므로(그러면 $M_{k-1}P_{k-1}\cdots M_1P_1A$의 앞 $k$열이 일차 종속이다) 피벗이 $0$이 아니고 $U$가 정칙이다.
\end{proof}

LAPACK은 $P$를 행렬로 들고 다니지 않고 정수 배열 \code{ipiv}로 기록한다. \code{ipiv(k)} $=p$는 ``$k$번째 단계에서 행 $k$와 행 $p$를 바꿨다''는 뜻이다(1부터 센다). 교환은 차례로 적용해야 뜻을 가지므로 \code{ipiv}는 순열 그 자체가 아니다. stardust의 \code{lu}는 이것을 ``$x$의 몇 번째 행이 $LU$의 $r$번째 행이 되는가''라는 순열 \code{perm}으로 바꿔 돌려준다(\cref{alg:lu-perm}).

\subsection{Growth Factor}

부분 피벗이 승수를 $1$ 이하로 묶는다고 해서 $U$의 원소까지 작게 묶이지는 않는다. 다음 절의 오차 해석에서 중요한 양은 소거 중에 나타나는 원소의 크기이며, 그 크기를 재는 것이 \term{성장 인자}{growth factor}
\[
\rho_n=\frac{\max_{i,j,k}|a^{(k)}_{ij}|}{\max_{i,j}|a_{ij}|}
\]
이다. $a^{(k)}_{ij}$는 $k$번째 단계 뒤의 원소이다.

\begin{proposition}
\label{prop:lu-growth-bound}
부분 피벗에서 $\rho_n\le2^{n-1}$이고 이 상한은 다음 행렬에서 등호가 된다.
\[
W_n=\begin{pmatrix}
1&&&&1\\
-1&1&&&1\\
-1&-1&1&&1\\
\vdots&&&\ddots&\vdots\\
-1&-1&\cdots&-1&1
\end{pmatrix}
\]
\end{proposition}

\begin{proof}
$|\ell_{ik}|\le1$이므로 $|a^{(k+1)}_{ij}|=|a^{(k)}_{ij}-\ell_{ik}a^{(k)}_{kj}|\le2\max|a^{(k)}|$이고, $n-1$번의 단계에서 최댓값이 매 단계 많아야 두 배가 된다. $W_n$에서는 $k$번째 단계에 $k$열의 대각 아래가 모두 $-1$이고 피벗이 $1$이다. 절댓값이 같은 후보 가운데 처음 것을 고르는 규칙(LAPACK의 \code{idamax})에서는 교환이 일어나지 않는다. 승수는 모두 $-1$이고 앞의 $n-1$열의 대각 아래가 $0$이 되는 동안 대각은 $1$로 남는다. 마지막 열은 $k$번째 단계에 대각 아래의 원소들이 모두 $a_{kn}^{(k)}$를 더해 받는다. 처음에 모두 $1$이므로 $k$번째 단계 뒤에 $a^{(k)}_{in}=2^{k}$($i>k$)이고 결국 $u_{nn}=2^{n-1}$이다.
\end{proof}

$n=64$이면 $2^{63}\approx9.2\times10^{18}$이고 binary64의 정밀도를 모두 잃는다. 그런데 실제 계산에서 이런 성장은 거의 보이지 않는다. \cref{fig:lu-growth}는 표준 정규 원소의 무작위 행렬에서 잰 성장이다. 소거 중 가장 큰 원소를 원소의 표준편차 $\sigma$로 나누면 중앙값이 $0.85\,n^{2/3}$ 곁에 놓인다(로그 축에서 맞춘 기울기는 $0.69$). Trefethen과 Schreiber가 관찰한 $n^{2/3}$ 법칙이다\cite{trefethen1990average}. 분모를 $\max|a_{ij}|$로 바꾼 $\rho_n$은 기울기가 $0.56$으로 조금 작은데, 정규 표본 $n^2$개의 최댓값이 $\sqrt{4\ln n}$처럼 자라 분모를 키우기 때문이다. 지수적 성장이 실제 응용에서 나타난 예(두 점 경계값 문제의 다중 사격법, Wright\cite{wright1993})가 있지만 드물고, 그 드묾을 설명하는 완결된 이론은 아직 없다. 이 간극이 부분 피벗이 ``이론상 불안정하지만 실제로는 안정한'' 알고리즘으로 불리는 까닭이다.

\begin{figure}[t]
\centering
\begin{tikzpicture}
\begin{loglogaxis}[width=0.86\textwidth,height=0.5\textwidth,
  xlabel={$n$},ylabel={성장},grid=major,grid style={gray!25},
  legend pos=south east,legend style={font=\small,draw=none,fill=white,fill opacity=0.85,text opacity=1},
  xmin=7,xmax=1400,ymin=1,ymax=200]
\addplot[only marks,mark=*,mark size=1.3pt,black] table[x=n,y=sigmed]{figures/48/growth.dat};
\addlegendentry{$\max|a^{(k)}_{ij}|/\sigma$의 중앙값}
\addplot[domain=7:1400,samples=40,dashed,black] {0.85*x^(2/3)};
\addlegendentry{$0.85\,n^{2/3}$}
\addplot[only marks,mark=square*,mark size=1.2pt,black!55] table[x=n,y=median]{figures/48/growth.dat};
\addlegendentry{$\rho_n$의 중앙값}
\addplot[only marks,mark=triangle*,mark size=1.5pt,black!55] table[x=n,y=max]{figures/48/growth.dat};
\addlegendentry{$\rho_n$의 최댓값}
\end{loglogaxis}
\end{tikzpicture}
\caption{표준 정규 원소($\sigma=1$)의 무작위 행렬에서 잰 부분 피벗의 성장. $n=10$에서 $2000$개, $n=1000$에서 $12$개의 행렬을 썼다(\code{figures/48/growth.py}). 상한 $2^{n-1}$은 $n=64$에서 이미 $10^{18}$을 넘지만 관찰된 성장은 수십 배에 그친다.}
\label{fig:lu-growth}
\end{figure}

부분 피벗 말고도 피벗 규칙이 있다. \term{완전 피벗}{complete pivoting}은 남은 부분행렬 전체에서 가장 큰 원소를 골라 행과 열을 함께 바꾼다. 성장 인자의 상한이 $n^{1/2+\ln n/4}$ 정도로 다항식에 가깝다는 Wilkinson의 결과가 있지만\cite{wilkinson1961error} 단계마다 $O((n-k)^2)$개의 원소를 훑어야 하므로 전체 비교 횟수가 $O(n^3)$이 되고, 그보다 나쁜 것은 블록 알고리즘으로 바꿀 수 없다는 점이다. 피벗을 고르려면 남은 행렬 전체가 최신이어야 하기 때문이다. 행과 열에서 번갈아 최대를 찾는 \term{룩 피벗}{rook pivoting}이 그 사이에 있다. LAPACK의 \code{getrf}와 그것을 부르는 stardust는 부분 피벗만 쓴다.

복소 행렬에는 한 가지가 더 있다. LAPACK의 \code{izamax}는 계산을 아끼려고 $|z|$ 대신 $|\operatorname{Re}z|+|\operatorname{Im}z|$가 가장 큰 원소를 고른다. 그러면 $|\ell_{ik}|\le1$이 아니라 $|\ell_{ik}|\le\sqrt2$만 보장된다. $|z|\le|\operatorname{Re}z|+|\operatorname{Im}z|\le\sqrt2|z|$이므로
\[
|\ell_{ik}|=\frac{|a_{ik}|}{|a_{kk}|}\le\frac{|\operatorname{Re}a_{ik}|+|\operatorname{Im}a_{ik}|}{(|\operatorname{Re}a_{kk}|+|\operatorname{Im}a_{kk}|)/\sqrt2}\le\sqrt2
\]
이기 때문이다. stardust의 시험은 복소 \code{lu}의 $L$을 이 $\sqrt2$로 확인한다.

\section{Backward Error Analysis}
\label{sec:lu-backward-error}

\subsection{The Error in the Factors}

11장의 표준 모형 $\mathrm{fl}(a\mathbin{\mathrm{op}}b)=(a\mathbin{\mathrm{op}}b)(1+\delta)$, $|\delta|\le u$와 기호 $\gamma_k=ku/(1-ku)$($ku<1$)를 쓴다. 소거와 대입에 나오는 계산은 모두 다음 꼴이다.

\begin{lemma}
\label{lem:lu-recurrence}
$y=(c-\sum_{i=1}^{k-1}a_ib_i)/b_k$를 $c$에서 차례로 $a_ib_i$를 빼고 마지막에 $b_k$로 나누어 계산하면 계산값 $\hat y$는
\[
\Bigl|c-\sum_{i=1}^{k-1}a_ib_i-b_k\hat y\Bigr|\le\gamma_k\Bigl(|b_k\hat y|+\sum_{i=1}^{k-1}|a_ib_i|\Bigr)
\]
을 만족한다. 나누지 않는 경우($b_k=1$)도 같다.
\end{lemma}

\begin{proof}
$\hat s_0=c$, $\hat s_i=\mathrm{fl}(\hat s_{i-1}-\mathrm{fl}(a_ib_i))$로 두면 $\hat s_i=(\hat s_{i-1}-a_ib_i(1+\epsilon_i))(1+\delta_i)$이고 $\hat y=\hat s_{k-1}(1+\delta_k)/b_k$이다. 점화식을 펼치면
\[
\hat s_{k-1}=c\prod_{j=1}^{k-1}(1+\delta_j)-\sum_{i=1}^{k-1}a_ib_i(1+\epsilon_i)\prod_{j=i}^{k-1}(1+\delta_j)
\]
이다. 양변을 $\prod_{j=1}^{k-1}(1+\delta_j)$로 나누고 $\hat s_{k-1}=b_k\hat y/(1+\delta_k)$를 넣으면
\[
b_k\hat y\,(1+\delta_k)^{-1}\prod_{j=1}^{k-1}(1+\delta_j)^{-1}=c-\sum_{i=1}^{k-1}a_ib_i\,(1+\epsilon_i)\prod_{j=1}^{i-1}(1+\delta_j)^{-1}
\]
이다. 왼쪽의 인수는 $(1+\delta)^{-1}$ 꼴 $k$개의 곱이고 오른쪽 $i$번째 항의 인수는 $(1+\delta)^{\pm1}$ 꼴 $i$개의 곱이다. 11장의 보조정리로 그런 곱은 $1+\theta$, $|\theta|\le\gamma_k$로 쓸 수 있으므로 $b_k\hat y(1+\theta_0)=c-\sum_ia_ib_i(1+\theta_i)$이고 $|\theta_i|\le\gamma_k$이다. 옮겨 쓰면 $c-\sum_ia_ib_i-b_k\hat y=b_k\hat y\,\theta_0+\sum_ia_ib_i\,\theta_i$이고 절댓값을 취하면 결론이다. 나눗셈이 없으면 $\delta_k$가 없어 $\gamma_{k-1}$로도 충분하다. 빼는 순서를 바꿔도 각 항이 겪는 반올림은 $k$번을 넘지 않으므로 같은 상한이 성립한다.
\end{proof}

LU의 원소는 $u_{kj}=a_{kj}-\sum_{i<k}\ell_{ki}u_{ij}$($k\le j$)와 $\ell_{ik}=(a_{ik}-\sum_{j<k}\ell_{ij}u_{jk})/u_{kk}$($i>k$)로 정해진다. 루프를 어떤 순서로 돌리든(\cref{sec:lu-unblocked}) 원소 하나는 이 꼴의 빼기와 한 번의 나눗셈으로 만들어지고, 빼는 순서는 보조정리의 상한에 영향을 주지 않으므로 보조정리가 그대로 적용된다.

\begin{theorem}[Wilkinson]
\label{thm:lu-backward}
피벗 없이 소거해 계산한 인수 $\hat L$, $\hat U$는
\[
\hat L\hat U=A+\Delta A,\qquad|\Delta A|\le\gamma_n|\hat L||\hat U|
\]
를 만족한다. 부분 피벗이면 $A$를 $PA$로 바꾸어 같은 식이 성립한다. 부등식과 절댓값은 원소별이다.
\end{theorem}

\begin{proof}
$k\le j$이면 \cref{lem:lu-recurrence}를 $c=a_{kj}$, $b_k=1$에 쓰면 $|a_{kj}-\sum_{i\le k}\hat\ell_{ki}\hat u_{ij}|\le\gamma_k\sum_{i\le k}|\hat\ell_{ki}||\hat u_{ij}|$이다($\hat\ell_{kk}=1$). $i>k$이면 $c=a_{ik}$, $b_k=\hat u_{kk}$에 써서 $|a_{ik}-\sum_{j\le k}\hat\ell_{ij}\hat u_{jk}|\le\gamma_k\sum_{j\le k}|\hat\ell_{ij}||\hat u_{jk}|$이다. 두 경우 모두 왼쪽 합은 $(\hat L\hat U)$의 해당 원소이고 오른쪽 합은 $(|\hat L||\hat U|)$의 해당 원소이며, $k\le n$이므로 $\gamma_k\le\gamma_n$이다. 부분 피벗은 소거 전에 행 순서를 정하는 것일 뿐 산술 연산을 바꾸지 않으므로 $PA$에 대해 같은 논증이 성립한다.
\end{proof}

\subsection{The Error in the Solution}

$Ax=b$는 $Ly=Pb$와 $Ux=y$ 두 번의 삼각 풀이로 푼다. 전진 대입 $y_i=(b_i-\sum_{j<i}\ell_{ij}y_j)/\ell_{ii}$에 \cref{lem:lu-recurrence}를 쓰면 같은 방법으로 계산값 $\hat y$가 $(L+\Delta L)\hat y=b$, $|\Delta L|\le\gamma_n|L|$을 만족한다. 후진 대입도 같다.

\begin{theorem}
\label{thm:lu-solve-backward}
부분 피벗 LU와 두 번의 대입으로 계산한 해 $\hat x$는
\[
(A+\Delta A)\hat x=b,\qquad|\Delta A|\le\gamma_{3n}P^{\mathsf T}|\hat L||\hat U|
\]
를 만족한다.
\end{theorem}

\begin{proof}
$P$를 $A$와 $b$에 미리 적용했다고 보고 $P=I$로 둔다. $(\hat L+\Delta L)\hat y=b$, $(\hat U+\Delta U)\hat x=\hat y$이므로
\[
b=(\hat L+\Delta L)(\hat U+\Delta U)\hat x=(A+\Delta A_1+\hat L\Delta U+\Delta L\hat U+\Delta L\Delta U)\hat x
\]
이고 $\Delta A_1=\hat L\hat U-A$이다. 원소별로 $|\Delta A_1|\le\gamma_n|\hat L||\hat U|$, $|\hat L\Delta U|+|\Delta L\hat U|\le2\gamma_n|\hat L||\hat U|$, $|\Delta L\Delta U|\le\gamma_n^2|\hat L||\hat U|$이다. $t=nu<1/3$으로 두면 $3\gamma_n+\gamma_n^2=(3t-2t^2)/(1-t)^2$이고, 이것이 $\gamma_{3n}=3t/(1-3t)$ 이하라는 것은 분모를 곱해 정리한 $3t(1-t)^2-(3t-2t^2)(1-3t)=5t^2-3t^3\ge0$에서 나온다.
\end{proof}

이 정리를 노름으로 바꾸면 성장 인자가 어디서 들어오는지 보인다. 부분 피벗이면 $|\hat\ell_{ij}|\le1$이고 $\hat U$의 원소는 소거 중의 원소이므로 $|\hat u_{ij}|\le\rho_n\max|a_{ij}|$이다. 그러면 $(|\hat L||\hat U|)_{ij}\le n\rho_n\max|a_{ij}|$이고 행의 합은 $n^2\rho_n\max|a_{ij}|\le n^2\rho_n\|A\|_\infty$ 이하이다. 곧
\begin{equation}
\label{eq:lu-normwise}
\frac{\|\Delta A\|_\infty}{\|A\|_\infty}\le n^2\rho_n\gamma_{3n}.
\end{equation}
$n^2$은 원소별 상한을 모두 같은 방향으로 더한 결과라 실제로는 훨씬 작고, 결국 관건은 $\rho_n$이다. 부분 피벗 LU가 실제로 안정한 것은 $\rho_n$이 실제로 작기 때문이다.

\subsection{Backward Error, Forward Error, and the Condition Number}

계산한 해의 \term{후진 오차}{backward error}는 그 해를 정확한 해로 만드는 가장 작은 섭동이다. 정규화한 정의
\[
\eta(\hat x)=\min\{\epsilon:(A+\Delta A)\hat x=b+\Delta b,\ \|\Delta A\|\le\epsilon\|A\|,\ \|\Delta b\|\le\epsilon\|b\|\}
\]
는 계산할 수 있는 식을 가진다.

\begin{theorem}[Rigal--Gaches]
\label{thm:lu-rigal-gaches}
벡터의 2-노름과 그에 따른 행렬 노름에서 $r=b-A\hat x$로 두면 $\eta(\hat x)=\|r\|/(\|A\|\|\hat x\|+\|b\|)$이다.
\end{theorem}

\begin{proof}
$(A+\Delta A)\hat x=b+\Delta b$이면 $r=\Delta A\hat x-\Delta b$이므로 $\|r\|\le\epsilon(\|A\|\|\hat x\|+\|b\|)$이고 $\eta$는 적어도 오른쪽 값이다. 반대로 $\epsilon_\ast=\|r\|/(\|A\|\|\hat x\|+\|b\|)$에 대해
\[
\Delta A=\frac{\|A\|\|\hat x\|}{\|A\|\|\hat x\|+\|b\|}\cdot\frac{r\hat x^{\mathsf H}}{\|\hat x\|^2},\qquad
\Delta b=-\frac{\|b\|}{\|A\|\|\hat x\|+\|b\|}\,r
\]
로 두면 $\Delta A\hat x-\Delta b=r$이어서 $(A+\Delta A)\hat x=b+\Delta b$이다. 계수 1인 $r\hat x^{\mathsf H}$의 2-노름은 $\|r\|\|\hat x\|$이므로 $\|\Delta A\|=\epsilon_\ast\|A\|$이고 $\|\Delta b\|=\epsilon_\ast\|b\|$이다. 그러므로 $\eta=\epsilon_\ast$이다.
\end{proof}

후진 오차가 작다고 해서 해가 정확한 것은 아니다. $(A+\Delta A)\hat x=b$이면 $\hat x-x=-A^{-1}\Delta A\hat x$이므로
\[
\frac{\|\hat x-x\|}{\|\hat x\|}\le\|A^{-1}\|\|\Delta A\|=\kappa(A)\frac{\|\Delta A\|}{\|A\|},\qquad\kappa(A)=\|A\|\|A^{-1}\|
\]
이다. 오른쪽의 $\hat x$를 $x$로 바꾼 꼴은 $\kappa\|\Delta A\|/\|A\|<1$일 때 분모에 $1-\kappa\|\Delta A\|/\|A\|$가 붙는다. 후진 안정한 알고리즘의 전진 오차는 조건수에 비례한다는 이 관계가 \cref{fig:lu-forward}에 그대로 나타난다. 후진 오차는 조건수와 상관없이 $u$ 근처에 머물고 전진 오차는 $\kappa(A)u$를 따라 오른다.

\begin{figure}[t]
\centering
\begin{tikzpicture}
\begin{loglogaxis}[width=0.86\textwidth,height=0.52\textwidth,
  xlabel={$\kappa_2(A)$},ylabel={상대오차},grid=major,grid style={gray!25},
  legend pos=north west,legend style={font=\small,draw=none,fill=none},
  ymin=1e-18,ymax=1]
\addplot[only marks,mark=*,mark size=1.1pt,black] table[x=kappa,y=fwd]{figures/48/errors.dat};
\addlegendentry{전진 오차 $\|\hat x-x\|_2/\|x\|_2$}
\addplot[only marks,mark=o,mark size=1.3pt,black!60] table[x=kappa,y=bwd]{figures/48/errors.dat};
\addlegendentry{후진 오차 $\eta(\hat x)$}
\addplot[domain=1:1e16,samples=20,dashed,black] {x*1.11e-16};
\addlegendentry{$\kappa_2(A)\,u$}
\end{loglogaxis}
\end{tikzpicture}
\caption{조건수를 정해 만든 $n=100$ 무작위 행렬에서 부분 피벗 LU로 푼 해의 오차. 만드는 방법은 \code{figures/48/errors.py}에 있다.}
\label{fig:lu-forward}
\end{figure}

\subsection{Iterative Refinement}

후진 오차가 이미 작다면 해를 더 정확하게 만드는 길은 조건수를 피하는 것밖에 없다. \term{반복 개선}{iterative refinement}은 이미 있는 LU 인수를 다시 써서 그것을 한다. 잔차 $r_i=b-Ax_i$를 구하고 $Ad_i=r_i$를 같은 인수로 풀어 $x_{i+1}=x_i+d_i$로 고친다. 계산한 $d_i$가 $(A+\Delta A_i)\hat d_i=r_i$를 만족한다고 하면
\[
x_{i+1}-x=x_i-x+(A+\Delta A_i)^{-1}A(x-x_i)=(A+\Delta A_i)^{-1}\Delta A_i\,(x_i-x)
\]
이다. 그러므로 잔차를 정확히 구한다면 오차는 한 번에 $\|A^{-1}\Delta A_i\|\lesssim\kappa(A)n^2\rho_n\gamma_{3n}$배로 줄어들고, 이 값이 $1$보다 충분히 작으면 몇 번 만에 수렴한다. 실제로 잔차를 작업 정밀도로 구하면 잔차의 반올림 오차가 한계를 정한다. 그래도 Skeel은 한 번의 개선만으로 원소별 후진 오차 $\omega(\hat x)=\min\{\epsilon:(A+\Delta A)\hat x=b+\Delta b,\ |\Delta A|\le\epsilon|A|,\ |\Delta b|\le\epsilon|b|\}$가 $O(u)$가 된다는 것을 보였다\cite{skeel1980refinement}. 증명의 줄거리는 이렇다. (i) 작업 정밀도로 구한 잔차는 $\hat r_0=r_0+\Delta r$, $|\Delta r|\le\gamma_{n+1}(|A||x_0|+|b|)$를 만족한다(\cref{lem:lu-recurrence}를 내적 하나에 쓴 것). (ii) 보정은 $(A+\Delta A)\hat d=\hat r_0$이고 $|\Delta A|\le\gamma_{3n}|\hat L||\hat U|$이다. (iii) 새 잔차는 $r_1=b-A(x_0+\hat d)=\hat r_0-\Delta r-A\hat d=\Delta A\,\hat d-\Delta r$이다. (iv) 보정 $\hat d$ 자체가 이미 처음 해의 오차만큼만 작으므로 $\Delta A\,\hat d$는 $u$의 제곱 크기이고, $|\hat L||\hat U|$와 $|A|$의 비와 $A^{-1}$의 원소별 조건수가 적당히 작다는 가정 아래 무시된다. 그러면 $|r_1|\lesssim|\Delta r|\le\gamma_{n+1}(|A||x_1|+|b|)$가 남고, 원소별 후진 오차의 Oettli--Prager 공식 $\omega(x)=\max_i|r_i|/(|A||x|+|b|)_i$로 $\omega(x_1)=O(u)$이다. (iv)의 가정을 정확히 적고 확인하는 일이 증명의 대부분이며, Higham의 책 12.2절이 이 논증을 처음부터 끝까지 따라간다\cite{higham2002accuracy}. 위 식이 보여 주는 핵심은 개선의 효과가 인수의 후진 오차 $\Delta A_i$를 한 번 더 곱해 줄이는 데서 온다는 것이다. LAPACK은 이것을 \code{gerfs}로 제공하지만 stardust의 \code{solve}는 개선을 하지 않는다. numpy와 scipy의 \code{solve}도 하지 않으므로 결과의 정밀도는 같은 수준이고, 개선이 필요한 사용자는 \code{lu}와 \code{lu\_solve}로 위의 반복을 직접 짤 수 있다.
\section{Unblocked Algorithms}
\label{sec:lu-unblocked}

\subsection{Six Loop Orders}

LU 분해의 산술은 세 겹의 루프 하나로 적힌다. 피벗을 잠시 잊으면
\[
\text{for } k,\ \text{for } i>k,\ \text{for } j>k:\quad a_{ij}\gets a_{ij}-\ell_{ik}a_{kj},\qquad \ell_{ik}=a_{ik}/a_{kk}
\]
이고 $(k,i,j)$의 여섯 가지 순서가 모두 같은 산술을 한다. 원소 $a_{ij}$가 받는 빼기는 언제나 $k=1,2,\dots$의 순서로 일어나므로 반올림까지 같은 결과가 나온다. 다른 것은 메모리를 읽고 쓰는 모양이다. \cref{fig:lu-looking}는 세 가지 대표 꼴을 보인다.
\begin{itemize}
\item \term{오른쪽 보기}{right-looking}($kij$, $kji$): $k$번째 단계에서 승수 열을 만들자마자 오른쪽 아래 남은 행렬 전체를 갱신한다. 갱신이 랭크 1 행렬 $\ell_ku_k^{\mathsf T}$를 빼는 것이어서 바깥곱 꼴이라고도 한다.
\item \term{왼쪽 보기}{left-looking}($jki$, $jik$): $j$열을 만들 차례가 되어서야 왼쪽에 이미 만든 열들로 $j$열 하나를 갱신한다. 한 단계에 쓰는 것은 한 열뿐이다.
\item \term{Crout 꼴}{Crout form}($ikj$, $ijk$ 등): $k$번째 단계에 $U$의 $k$행과 $L$의 $k$열을 이미 만든 부분과의 내적으로 만든다.
\end{itemize}

\begin{figure}[t]
\centering
\begin{tikzpicture}[x=0.36cm,y=0.36cm,font=\small]
\foreach \xo/\name in {0/오른쪽 보기,9/왼쪽 보기,18/Crout 꼴}{
  \draw (\xo,0) rectangle (\xo+7,7);
  \node at (\xo+3.5,-1) {\name};
}
% right-looking: done L cols + U rows in light gray; current col/row dark; trailing updated hatched
\fill[gray!15] (0,0) rectangle (3,7); \fill[gray!15] (3,4) rectangle (7,7);
\fill[gray!55] (3,0) rectangle (4,4); \fill[gray!55] (3,3) rectangle (7,4);
\fill[pattern=north east lines] (4,0) rectangle (7,3); \draw (4,0) rectangle (7,3);
% left-looking: done cols light; read L part of done cols; column j written
\fill[gray!15] (9,0) rectangle (12,7);
\fill[gray!55] (12,0) rectangle (13,7);
\draw[pattern=north east lines] (12,0) rectangle (13,7);
% crout: row k of U and column k of L written; done rows/cols read
\fill[gray!15] (18,0) rectangle (21,7); \fill[gray!15] (21,4) rectangle (25,7);
\fill[gray!55] (21,0) rectangle (22,4); \fill[gray!55] (21,3) rectangle (25,4);
\end{tikzpicture}
\caption{세 가지 루프 꼴이 한 단계에 읽고 쓰는 곳. 연한 회색은 이미 끝난 $L$과 $U$, 진한 회색은 이 단계에 완성되는 부분, 빗금은 이 단계에 갱신되는 부분이다. 오른쪽 보기는 남은 행렬 전체를 매 단계 고치고, 왼쪽 보기는 한 열만 고치며, Crout 꼴은 남은 행렬을 건드리지 않는다.}
\label{fig:lu-looking}
\end{figure}

열 우선(column-major) 배열에서는 같은 열의 원소가 메모리에 붙어 있으므로 가장 안쪽 루프가 $i$(행)를 돌아야 연속한 메모리를 훑는다. 그래서 LAPACK의 꼴은 $kji$(오른쪽 보기)와 $jki$(왼쪽 보기)이다. 부분 피벗을 넣으면 다음과 같다.

\begin{algorithm}[t]
\begin{algorithmic}[1]
\Require $m\times n$ 행렬 $A$ (열 우선, 그 자리에 덮어씀)
\Ensure $A$의 대각 아래에 $L$의 승수, 대각과 위에 $U$; 교환 기록 $\mathit{ipiv}$; 정확히 $0$인 첫 피벗의 번호 $\mathit{info}$
\State $\mathit{info}\gets0$
\For{$k=1,\dots,\min(m,n)$}
  \State $p\gets\arg\max_{k\le i\le m}|a_{ik}|$ \Comment{같은 값이면 가장 앞의 $i$}
  \State $\mathit{ipiv}_k\gets p$
  \If{$a_{pk}=0$}
    \State \textbf{if} $\mathit{info}=0$ \textbf{then} $\mathit{info}\gets k$ \Comment{특이: 이 열은 소거하지 않는다}
    \State \textbf{continue}
  \EndIf
  \State 행 $k$와 행 $p$를 열 $1,\dots,n$ 전체에서 바꾼다
  \State $a_{k+1:m,\,k}\gets a_{k+1:m,\,k}/a_{kk}$ \Comment{승수}
  \For{$j=k+1,\dots,n$} \Comment{랭크 1 갱신, 열 하나씩}
    \For{$i=k+1,\dots,m$}
      \State $a_{ij}\gets a_{ij}-a_{ik}\,a_{kj}$
    \EndFor
  \EndFor
\EndFor
\end{algorithmic}
\caption{부분 피벗을 쓰는 오른쪽 보기 LU ($kji$ 꼴, LAPACK \code{getf2}의 산술)}
\label{alg:lu-right}
\end{algorithm}

\cref{alg:lu-right}의 줄마다 할 일을 짚는다.
\begin{itemize}
\item 피벗 찾기(3행)는 $k$열의 대각부터 아래를 훑는다. 절댓값이 같은 후보가 여럿이면 처음 것을 고르는 규칙은 결과를 정하는 규칙이기도 하다. \cref{prop:lu-growth-bound}의 $W_n$이 교환 없이 소거되는 것이 이 규칙 때문이다.
\item 피벗이 정확히 $0$이면(5행) 그 열의 대각 아래도 모두 $0$이므로 소거할 것이 없다. 알고리즘은 멈추지 않고 첫 $0$ 피벗의 번호만 기록한다. 그래서 특이 행렬도 \cref{sec:lu-factorization}의 뜻에서 분해가 끝까지 나오고, 그것을 에러로 볼지는 부르는 쪽이 정한다.
\item 행 교환(8행)은 이미 끝난 왼쪽 열들(승수)까지 바꾼다. \cref{thm:lu-partial-pivoting}의 증명에서 승수를 뒤의 교환으로 섞어야 했던 것이 바로 이 일이다. 이 덕분에 끝난 뒤 $A$의 대각 아래가 그대로 $PA=LU$의 $L$이 된다.
\item 승수(9행)는 피벗으로 나눈다. LAPACK의 \code{getf2}는 $1/a_{kk}$를 한 번 구해 곱하는데, $|a_{kk}|$가 가장 작은 정규 수보다 작으면 역수가 넘치므로 그때만 원소마다 나눈다. 곱셈이 나눗셈보다 훨씬 빠르기 때문에 고른 길이며, 역수를 곱한 결과는 나눈 결과와 반올림 한 번만큼 다를 수 있다.
\item 갱신(10--14행)은 열마다 \code{axpy} 한 번이다. LAPACK은 이 이중 루프를 BLAS의 랭크 1 갱신 \code{ger} 하나로 부른다.
\end{itemize}

\subsection{Operation Count and Memory Traffic}

$k$번째 단계는 나눗셈 $m-k$번과 갱신 $(m-k)(n-k)$개 원소에 곱셈과 뺄셈 하나씩이다. 정사각이면
\[
\sum_{k=1}^{n-1}\bigl[(n-k)+2(n-k)^2\bigr]=\frac{n(n-1)}2+\frac{(n-1)n(2n-1)}3=\frac23n^3-\frac12n^2-\frac16n
\]
플롭이다. 이 $(2/3)n^3$이 LU의 비용으로 늘 인용되는 수이다.

산술보다 중요한 것이 메모리이다. 오른쪽 보기의 $k$번째 단계는 남은 $(n-k)^2$개 원소를 모두 읽고 쓰는데 원소 하나에 플롭 둘이다. 남은 행렬이 캐시보다 크면 이 갱신은 메모리 대역폭에 묶이고, 산술 장치는 데이터를 기다리며 논다. \cref{fig:lu-gflops}의 맨 아래 곡선이 그것이다. 직접 짠 \cref{alg:lu-right}은 $n$이 커져도 초당 $12\times10^9$ 플롭 남짓에서 멈추고, 같은 산술을 하는 LAPACK의 \code{dgetf2}도 같은 곳에 있다. 이 상자의 처리기는 한 코어로도 그 열 배 넘게 계산할 수 있다. 왼쪽 보기는 쓰는 양이 적은 대신 매 단계 왼쪽의 모든 열을 읽으므로 메모리 이동의 총량은 같은 차수이다. 이 벽을 넘는 길은 루프 순서가 아니라 블록이다.

\section{Blocked and Recursive LU}
\label{sec:lu-blocked}

\subsection{The BLAS Levels}

BLAS(Basic Linear Algebra Subprograms)는 선형대수의 기본 연산을 셋으로 나눈다\cite{lawson1979blas1,dongarra1990blas3}. 1단계는 벡터 연산(\code{axpy}, \code{dot}; $O(n)$ 데이터에 $O(n)$ 연산), 2단계는 행렬–벡터 연산(\code{gemv}, \code{ger}, \code{trsv}; $O(n^2)$ 데이터에 $O(n^2)$ 연산), 3단계는 행렬–행렬 연산(\code{gemm}, \code{trsm}; $O(n^2)$ 데이터에 $O(n^3)$ 연산)이다. 데이터 하나를 메모리에서 가져와 몇 번의 연산에 쓰는지(산술 강도)가 1, 2단계에서는 상수이고 3단계에서만 $n$에 비례한다. 3단계 연산은 행렬을 캐시와 레지스터에 맞는 블록으로 잘라 블록 하나를 여러 번 재사용할 수 있으므로 산술 장치의 최고 속도에 가깝게 돈다. 빠른 LU란 산술의 대부분을 \code{gemm}으로 바꾼 LU이다.

\subsection{Block Right-Looking LU}

$A$의 앞 $b$열을 \term{패널}{panel}로 떼어 블록으로 나눈다.
\[
A=\begin{pmatrix}A_{11}&A_{12}\\A_{21}&A_{22}\end{pmatrix},\qquad A_{11}\in\mathbb C^{b\times b}.
\]
패널 $\binom{A_{11}}{A_{21}}$을 부분 피벗으로 분해해 $P_1\binom{A_{11}}{A_{21}}=\binom{L_{11}}{L_{21}}U_{11}$을 얻었다고 하자. 같은 교환을 오른쪽 열에도 적용한 뒤 $P_1A=LU$의 블록 식을 쓰면
\[
\begin{pmatrix}A_{11}&A_{12}\\A_{21}&A_{22}\end{pmatrix}
=\begin{pmatrix}L_{11}&0\\L_{21}&I\end{pmatrix}
\begin{pmatrix}U_{11}&U_{12}\\0&S\end{pmatrix},\qquad
U_{12}=L_{11}^{-1}A_{12},\quad S=A_{22}-L_{21}U_{12}
\]
이다. 식 \eqref{eq:lu-schur-step}을 블록으로 바꾼 것이며 남은 일은 $S$를 같은 방법으로 분해하는 것이다(\cref{fig:lu-blocked}). $U_{12}$를 구하는 것은 단위 하삼각 행렬로 여러 우변을 푸는 \code{trsm}이고, $S$의 갱신은 $(n-b)\times b$와 $b\times(n-b)$의 곱인 \code{gemm}이다.

\begin{figure}[t]
\centering
\begin{tikzpicture}[x=0.42cm,y=0.42cm,font=\small]
\draw (0,0) rectangle (9,9);
\fill[gray!15] (0,0) rectangle (2,9); \fill[gray!15] (2,7) rectangle (9,9);
\fill[gray!55] (2,0) rectangle (4,7);
\fill[gray!35] (4,5) rectangle (9,7);
\fill[pattern=north east lines] (4,0) rectangle (9,5); \draw (4,0) rectangle (9,5);
\draw (2,0) -- (2,9) (4,0) -- (4,7) (2,7) -- (9,7) (2,5) -- (9,5);
\node at (1,4.5) {끝남};
\node[fill=white,inner sep=1pt] at (3,6) {$A_{11}$};
\node[fill=white,inner sep=1pt] at (3,2.5) {$A_{21}$};
\node at (6.5,6) {$A_{12}$};
\node[fill=white,inner sep=1pt] at (6.5,2.5) {$A_{22}$};
\begin{scope}[xshift=4.6cm]
\node[anchor=west,align=left,font=\small] at (0,6.8) {(i) 패널 분해 \code{getrf2}: 진한 회색};
\node[anchor=west,align=left,font=\small] at (0,5.3) {(ii) 교환을 좌우 열에 \code{laswp}};
\node[anchor=west,align=left,font=\small] at (0,3.8) {(iii) $U_{12}=L_{11}^{-1}A_{12}$ \code{trsm}: 중간 회색};
\node[anchor=west,align=left,font=\small] at (0,2.3) {(iv) $A_{22}\gets A_{22}-L_{21}U_{12}$ \code{gemm}: 빗금};
\end{scope}
\end{tikzpicture}
\caption{블록 오른쪽 보기 LU의 한 단계. 패널은 폭 $b$의 열 묶음이고 이 단계의 산술 대부분은 빗금 친 부분의 \code{gemm}이다.}
\label{fig:lu-blocked}
\end{figure}

\begin{algorithm}[t]
\begin{algorithmic}[1]
\Require $m\times n$ 행렬 $A$, 블록 폭 $b$ (LAPACK 참조 구현의 기본값 $64$)
\Ensure \cref{alg:lu-right}과 같은 꼴의 $L$, $U$, $\mathit{ipiv}$, $\mathit{info}$
\If{$b\le1$ \textbf{or} $b\ge\min(m,n)$}
  \State \Return \Call{getrf2}{$A$} \Comment{작으면 블록 없이}
\EndIf
\For{$j=1,\,1+b,\,1+2b,\dots\le\min(m,n)$}
  \State $j_b\gets\min(\min(m,n)-j+1,\ b)$
  \State \Call{getrf2}{$A_{j:m,\,j:j+j_b-1}$} \Comment{패널; 피벗 번호는 패널 안에서 센다}
  \State $\mathit{ipiv}_{j:j+j_b-1}\gets\mathit{ipiv}_{j:j+j_b-1}+j-1$; $\mathit{info}$도 같이 옮긴다
  \State 행 교환 $\mathit{ipiv}_{j:j+j_b-1}$을 열 $1,\dots,j-1$에 적용 \Comment{\code{laswp}: 끝난 $L$}
  \If{$j+j_b\le n$}
    \State 같은 교환을 열 $j+j_b,\dots,n$에 적용 \Comment{\code{laswp}}
    \State $A_{j:j+j_b-1,\,j+j_b:n}\gets L_{11}^{-1}A_{j:j+j_b-1,\,j+j_b:n}$ \Comment{\code{trsm}, $L_{11}$은 단위 하삼각}
    \If{$j+j_b\le m$}
      \State $A_{j+j_b:m,\,j+j_b:n}\gets A_{j+j_b:m,\,j+j_b:n}-A_{j+j_b:m,\,j:j+j_b-1}A_{j:j+j_b-1,\,j+j_b:n}$ \Comment{\code{gemm}}
    \EndIf
  \EndIf
\EndFor
\end{algorithmic}
\caption{블록 오른쪽 보기 LU (LAPACK \code{dgetrf}의 구조를 그대로 옮김)}
\label{alg:lu-blocked}
\end{algorithm}

\cref{alg:lu-blocked}은 LAPACK 참조 구현 \code{dgetrf}의 순서를 그대로 옮긴 것이다(\cite{anderson1999lapack}, 참조 구현의 소스로 확인). 몇 가지를 짚는다.
\begin{itemize}
\item 패널 분해는 $(m-j+1)\times j_b$의 길쭉한 행렬에 대한 LU이고 피벗을 패널 높이 전체에서 찾는다. 그래서 블록 알고리즘도 피벗 규칙은 \cref{alg:lu-right}과 같은 부분 피벗이며, 반올림을 빼면 같은 $P$, $L$, $U$를 낸다.
\item 패널이 돌려준 피벗 번호는 패널의 첫 행을 $1$로 센 것이므로 7행에서 $j-1$을 더해 전체 행 번호로 바꾼다.
\item 교환은 두 번에 나누어 적용한다. 패널 안의 열은 패널 분해가 이미 바꿨고, 왼쪽의 끝난 열(8행)과 오른쪽의 아직 손대지 않은 열(10행)을 따로 바꾼다.
\item 산술량을 세 부분으로 나누면 패널 분해가 단계마다 약 $(n-j)j_b^2$, \code{trsm}이 약 $j_b^2(n-j)$이고 나머지가 모두 \code{gemm}이다. 단계들을 더하면 \code{gemm}이 아닌 부분이 약 $n^2b$이므로 \code{gemm}의 몫은
\[
1-\frac{n^2b}{2n^3/3}=1-\frac{3b}{2n}
\]
이다. $n=1024$, $b=64$이면 $91\%$이다. $b$를 키우면 이 몫이 줄고, 줄이면 \code{gemm} 하나하나가 납작해져(내적 길이 $b$) 재사용이 줄어든다. $b$는 이 둘 사이의 타협이고 하드웨어마다 다르다. 참조 구현은 \code{ilaenv}가 돌려주는 $64$를 쓴다.
\end{itemize}

\cref{fig:lu-gflops}는 이 차이를 보인다. 직접 짠 \cref{alg:lu-blocked}(패널은 \cref{alg:lu-right}, \code{trsm}과 \code{gemm}은 Accelerate)이 블록 없는 루프보다 $n=1536$에서 일곱 배 빠르다. 같은 그림에서 $n=256$, $512$, $1024$의 블록 곡선이 움푹 꺼지는 것은 우연이 아니다. 이때 열 하나의 길이가 $2^{k}$의 배수여서 같은 행의 원소들이 캐시의 같은 묶음(set)에 몰린다. leading dimension만 $512$에서 $520$으로 늘려 같은 분해를 하면 초당 $36\times10^9$에서 $85\times10^9$ 플롭으로, $1024$에서 $1032$로 늘리면 $59\times10^9$에서 $122\times10^9$ 플롭으로 뛴다. 같은 산술, 같은 메모리 양인데 주소의 모양만 바뀐 결과이다.

\begin{figure}[t]
\centering
\begin{tikzpicture}
\begin{semilogxaxis}[width=0.86\textwidth,height=0.55\textwidth,
  xlabel={$n$},ylabel={$10^9$ 플롭/초},grid=major,grid style={gray!25},
  legend pos=north west,legend style={font=\small,draw=none,fill=none},
  xmin=28,xmax=2300,ymin=0,ymax=300,log basis x=2]
\addplot[black,mark=*,mark size=1.2pt] table[x=n,y=getrf]{figures/48/flops.dat};
\addlegendentry{Accelerate \code{dgetrf}}
\addplot[black,dashed,mark=square*,mark size=1.2pt] table[x=n,y=blocked]{figures/48/flops.dat};
\addlegendentry{\cref{alg:lu-blocked} ($b=64$)}
\addplot[black!60,mark=triangle*,mark size=1.4pt] table[x=n,y=getf2]{figures/48/flops.dat};
\addlegendentry{Accelerate \code{dgetf2}}
\addplot[black!60,dotted,thick,mark=o,mark size=1.2pt] table[x=n,y=unblocked]{figures/48/flops.dat};
\addlegendentry{\cref{alg:lu-right} (C 루프)}
\end{semilogxaxis}
\end{tikzpicture}
\caption{정사각 LU의 속도(Apple M4 Pro, 플롭 수는 $(2/3)n^3$으로 셈). 측정 코드는 \code{figures/48/bench.c}이다. \code{dgetrf}는 $n\ge768$에서 Accelerate가 안에서 여러 코어를 쓴다.}
\label{fig:lu-gflops}
\end{figure}

\subsection{Recursive LU}

블록 알고리즘의 패널 분해는 여전히 2단계 BLAS이다. 패널은 높이가 $n$이고 폭이 $b$이므로 패널의 산술이 전체의 $3b/(2n)$밖에 안 돼도 메모리에 묶인 속도로 돌기 때문에 시간으로는 그보다 큰 몫을 차지한다. Toledo와 Gustavson은 패널을 다시 반으로 나누는 재귀로 이것을 없앴다\cite{toledo1997recursive}. 열을 반으로 나누어 왼쪽 반을 재귀로 분해하고, 오른쪽 반을 \code{trsm}과 \code{gemm}으로 갱신한 뒤, 남은 아래 오른쪽을 재귀로 분해한다. LAPACK 3.6부터 패널 분해 \code{getrf2}가 이 재귀이다.

\begin{algorithm}[t]
\begin{algorithmic}[1]
\Function{getrf2}{$A\in\mathbb C^{m\times n}$}
  \If{$m=1$}
    \State $\mathit{ipiv}_1\gets1$; \textbf{if} $a_{11}=0$ \textbf{then} $\mathit{info}\gets1$
  \ElsIf{$n=1$} \Comment{한 열: 피벗과 나눗셈뿐}
    \State $p\gets\arg\max_i|a_{i1}|$; $\mathit{ipiv}_1\gets p$
    \If{$a_{p1}\ne0$}
      \State $a_{11}\leftrightarrow a_{p1}$
      \State \textbf{if} $|a_{11}|\ge\mathrm{sfmin}$ \textbf{then} $a_{2:m,1}\gets a_{2:m,1}\cdot(1/a_{11})$ \textbf{else} $a_{2:m,1}\gets a_{2:m,1}/a_{11}$
    \Else\ $\mathit{info}\gets1$
    \EndIf
  \Else
    \State $n_1\gets\lfloor\min(m,n)/2\rfloor$; $n_2\gets n-n_1$
    \State \Call{getrf2}{$A_{:,1:n_1}$} \Comment{왼쪽 반 $\binom{A_{11}}{A_{21}}$}
    \State 교환 $\mathit{ipiv}_{1:n_1}$을 $A_{:,n_1+1:n}$에 적용 \Comment{\code{laswp}}
    \State $A_{12}\gets L_{11}^{-1}A_{12}$ \Comment{\code{trsm}}
    \State $A_{22}\gets A_{22}-A_{21}A_{12}$ \Comment{\code{gemm}}
    \State \Call{getrf2}{$A_{22}$}; $\mathit{ipiv}_{n_1+1:\min(m,n)}\mathrel{+}=n_1$ \Comment{$\mathit{info}$도 $n_1$만큼 옮긴다}
    \State 교환 $\mathit{ipiv}_{n_1+1:\min(m,n)}$을 $A_{:,1:n_1}$에 적용 \Comment{\code{laswp}: 왼쪽의 $L_{21}$}
  \EndIf
\EndFunction
\end{algorithmic}
\caption{재귀 LU (LAPACK \code{dgetrf2})}
\label{alg:lu-recursive}
\end{algorithm}

\cref{alg:lu-recursive}의 기저 경우는 둘이다. 한 행($m=1$)이면 할 일이 없고, 한 열($n=1$)이면 피벗을 찾아 바꾸고 나누는 것이 전부이다. 여기서 $\mathrm{sfmin}$은 역수가 넘치지 않는 가장 작은 수(LAPACK의 \code{dlamch('S')})이고 \cref{alg:lu-right}의 승수 계산과 같은 이유로 둔다. 재귀의 모든 층에서 산술의 대부분이 \code{gemm}이고, 크기가 반씩 줄어드므로 어느 층의 부분행렬은 반드시 캐시에 맞는다. 그래서 블록 폭 $b$를 정할 필요가 없고(이런 알고리즘을 캐시를 모르는(cache-oblivious) 알고리즘이라 한다) 반올림을 빼면 결과는 다시 부분 피벗의 $P$, $L$, $U$이다. 나누는 기준이 $\min(m,n)/2$인 것은 길쭉한 패널($m\gg n$)에서도 열을 반으로 나누기 위해서이다.

\subsection{What Accelerate Does}

stardust는 LU를 직접 짜지 않고 Apple Accelerate의 LAPACK을 부른다. Accelerate의 \code{getrf}는 참조 구현과 같은 인터페이스와 같은 피벗 규칙(\code{ipiv}, 같은 값이면 앞의 행)을 가지지만 내부의 블록 크기, 재귀의 모양, 행렬 곱의 커널은 공개되어 있지 않다. 측정으로 알 수 있는 것은 이렇다. \cref{fig:lu-gflops}에서 \code{dgetrf}는 $n=1536$에 초당 $2.8\times10^{11}$ 플롭으로 직접 짠 블록 LU의 세 배이고, $n\ge768$에서 여러 코어를 쓴다. 블록 없는 \code{dgetf2}는 직접 짠 루프와 같은 메모리의 벽에 있다. 그리고 다음 절에서 보듯 작은 행렬에서 \code{getrf}는 이상한 행동을 한다. 그 행동이 stardust의 구현에 규칙 두 개를 더했다(\cref{sec:lu-batches}).

\section{Using the Factors}
\label{sec:lu-using}

\subsection{Triangular Solves and Row Interchanges}

$PA=LU$를 얻으면 $Ax=b$는 $Ly=Pb$, $Ux=y$이다. 전진 대입 $y_i=(Pb)_i-\sum_{j<i}\ell_{ij}y_j$와 후진 대입 $x_i=(y_i-\sum_{j>i}u_{ij}x_j)/u_{ii}$는 각각 $n^2$ 플롭이다(나눗셈 $n$번 포함). 우변이 $k$개면 $2kn^2$이고, 분해의 $(2/3)n^3$보다 작으므로 분해를 한 번 하고 여러 우변에 다시 쓰는 것이 이득이다. LAPACK의 \code{getrs}가 이 일을 한다.

\begin{algorithm}[t]
\begin{algorithmic}[1]
\Require \code{getrf}의 결과($L$, $U$를 담은 $A$와 $\mathit{ipiv}$), $n\times k$ 우변 $B$
\Ensure $B$에 $X=A^{-1}B$를 덮어씀
\For{$i=1,\dots,n$} \Comment{\code{laswp}: 교환을 기록된 순서대로}
  \State \textbf{if} $\mathit{ipiv}_i\ne i$ \textbf{then} $B$의 행 $i$와 행 $\mathit{ipiv}_i$를 바꾼다
\EndFor
\For{$j=1,\dots,k$} \Comment{열 우선: 우변 하나씩}
  \For{$i=1,\dots,n$} \Comment{$Ly=Pb$, $\ell_{ii}=1$이라 나누지 않는다}
    \State $b_{i+1:n,j}\gets b_{i+1:n,j}-b_{ij}\,a_{i+1:n,\,i}$
  \EndFor
  \For{$i=n,\dots,1$} \Comment{$Ux=y$}
    \State $b_{ij}\gets b_{ij}/a_{ii}$; \ $b_{1:i-1,j}\gets b_{1:i-1,j}-b_{ij}\,a_{1:i-1,\,i}$
  \EndFor
\EndFor
\end{algorithmic}
\caption{인수로 풀기 (\code{getrs}의 \code{'N'} 경우; 열 지향의 대입)}
\label{alg:lu-solve}
\end{algorithm}

\cref{alg:lu-solve}의 대입은 열 지향으로 적었다. $y_i$가 정해지는 순간 그 값을 $L$의 $i$열에 곱해 아래 원소들에서 한꺼번에 빼므로 안쪽 루프가 열 우선 배열의 연속한 메모리를 지난다. 산술은 행 지향($y_i=b_i-\sum_j\ell_{ij}y_j$)과 같고 빼는 순서만 다르다. 행 교환은 \code{ipiv}를 앞에서부터 차례로 적용해야 한다. 앞의 교환이 뒤의 교환이 가리키는 행의 내용을 바꾸기 때문이다.

교환 목록을 순열로 바꾸는 것은 단위 순열에 같은 교환을 차례로 적용하는 것이다.

\begin{algorithm}[t]
\begin{algorithmic}[1]
\Require 교환 목록 $\mathit{ipiv}_1,\dots,\mathit{ipiv}_p$ (1부터), 행 수 $m$
\Ensure $x_{\mathit{perm}_r,:}=(LU)_{r,:}$인 순열 $\mathit{perm}$ (0부터)
\State $\mathit{perm}_r\gets r$ ($r=0,\dots,m-1$)
\For{$r=0,\dots,p-1$}
  \State $\mathit{perm}_r\leftrightarrow\mathit{perm}_{\mathit{ipiv}_{r+1}-1}$
\EndFor
\end{algorithmic}
\caption{교환 목록에서 순열로 (stardust \code{lu}의 \code{perm})}
\label{alg:lu-perm}
\end{algorithm}

\cref{alg:lu-perm}이 맞는 까닭은 $PA$의 $r$번째 행이 원래 몇 번째 행이었는지를 따라가는 것이 바로 $A$의 행 번호들에 같은 교환을 적용하는 것이기 때문이다. 교환 $r\leftrightarrow\mathit{ipiv}_r$은 행렬의 행 내용을 바꾸고 동시에 ``지금 $r$번째 자리에 있는 원래 행의 번호''도 바꾼다. 결과는 numpy·scipy 관례의 \code{x[perm] = L @ U}이다.

\subsection{The Inverse}

역행렬은 $AX=I$를 푸는 것이다. 우변이 $n$개이므로 \cref{alg:lu-solve}을 그대로 쓰면 풀이에 $2n^3$ 플롭, 분해까지 $(8/3)n^3$이다. LAPACK의 \code{getri}는 다른 길을 간다. $U^{-1}$을 그 자리에서 구하고(\code{trtri}, $(1/3)n^3$), $XL=U^{-1}$을 오른쪽에서 풀고($n^3$), 마지막에 열을 \code{ipiv}의 역순으로 바꾼다. 그러면 분해를 더해 $2n^3$이고, 단위행렬의 $0$을 연산하지 않으니 풀이 쪽이 약 $(4/3)n^3$만큼 싸다. 반면 \code{getri}는 단계가 셋이어서 작은 행렬에서는 호출 비용이 크다. \cref{tab:lu-inverse}은 두 길을 직접 잰 것이다. $n\le1024$에서는 단위행렬을 우변으로 둔 풀이가 빠르거나 같고 $n=2048$에서 \code{getri}가 앞선다. stardust가 어느 쪽을 골랐고 왜 그랬는지는 \cref{sec:lu-stardust-family}에 있다.

\begin{table}[t]
\begin{sdtable}{colspec={rrrr}, row{1}={font=\small\bfseries}}
\toprule
$n$ & \code{getrf} & \code{getri} & \code{getrs}$(I)$ \\
\midrule
$8$ & $0.50$ & $1.21$ & $0.21$ \\
$64$ & $13.3$ & $41.4$ & $11.6$ \\
$256$ & $125$ & $187$ & $179$ \\
$1024$ & $2769$ & $4954$ & $3825$ \\
$2048$ & $14445$ & $34501$ & $38641$ \\
\bottomrule
\end{sdtable}
\caption{역행렬의 두 길(Accelerate, 한 번의 호출 시간, 마이크로초, 여러 번 잰 최솟값). \code{getri}와 \code{getrs}$(I)$는 같은 \code{getrf} 결과에서 시작한다(\code{figures/48/bench.c}의 \code{inv}).}
\label{tab:lu-inverse}
\end{table}

\subsection{Determinants}

$PA=LU$에서 $\det P=(-1)^{s}$($s$는 실제로 일어난 교환의 수, 곧 $\mathit{ipiv}_i\ne i$인 $i$의 수)이고 $\det L=1$이므로
\[
\det A=(-1)^{s}\prod_{i=1}^nu_{ii}
\]
이다. 곱을 그대로 계산하면 쉽게 넘치거나 언더플로한다. $\operatorname{diag}(10^{200},10^{200},10^{-300})$의 행렬식은 $10^{100}$으로 binary64 안에 있지만 앞의 두 원소를 곱하는 순간 $10^{400}$이 넘쳐 $\infty$가 된다. 크기가 $1000$인 행렬에서 대각이 모두 $0.4$ 근처면 행렬식은 $10^{-400}$ 부근이어서 정말로 언더플로한다. 그래서 두 가지를 구분한다. 행렬식 자체가 범위를 벗어나는 경우와 중간 곱만 벗어나는 경우이다. 뒤의 경우를 정확히 처리하려면 지수를 따로 들고 다닌다.

\begin{algorithm}[t]
\begin{algorithmic}[1]
\Require $U$의 대각 $u_1,\dots,u_n$, 교환 목록 $\mathit{ipiv}$, 특이 여부
\Ensure $\det A$ (중간 곱이 범위를 벗어나도 결과가 범위 안이면 그 값)
\If{특이} \Return $0$ \EndIf
\State $d\gets1$; \textit{ok} $\gets$ \textbf{true}
\For{$i=1,\dots,n$}
  \State $d\gets d\cdot u_i$; \textbf{if} $\mathit{ipiv}_i\ne i$ \textbf{then} $d\gets-d$
  \State \textit{ok} $\gets$ \textit{ok} \textbf{and} ($d$가 $0$이 아닌 정규 수)
\EndFor
\If{\textit{ok}} \Return $d$ \EndIf \Comment{대부분은 여기서 끝난다}
\State $d\gets1$; $e\gets0$ \Comment{$\det A=d\cdot2^{e}$로 들고 다닌다}
\For{$i=1,\dots,n$}
  \State $u_i=m_i\cdot2^{k_i}$, $|m_i|\in[1/2,1)$로 나누고 $e\gets e+k_i$ \Comment{\code{frexp}: 정확}
  \State $d\gets d\cdot m_i$; $d=m\cdot2^{k}$로 다시 나누어 $d\gets m$, $e\gets e+k$
  \State \textbf{if} $\mathit{ipiv}_i\ne i$ \textbf{then} $d\gets-d$
\EndFor
\State \Return $d\cdot2^{\max(-4000,\min(e,4000))}$ \Comment{\code{ldexp}: 여기서만 넘치거나 언더플로}
\end{algorithmic}
\caption{넘침을 피하는 행렬식 (stardust \code{det}의 \code{det\_of})}
\label{alg:lu-det}
\end{algorithm}

\cref{alg:lu-det}의 두 번째 루프는 첫 번째 루프와 정확히 같은 반올림을 한다. 2의 거듭제곱을 곱하거나 나누는 것은 결과가 정규 범위에 있는 한 정확하고, $|m_i|,|d|\in[1/2,1)$이면 곱도 $[1/4,1)$에 있어 정규 범위를 벗어나지 않기 때문이다. 그래서 첫 루프가 끝까지 정규 수였던 행렬에서는 두 루프의 가수가 비트까지 같고, 첫 루프가 중간에 넘쳤던 행렬에서는 두 번째 루프가 넘치지 않았다면 얻었을 값을 준다. 마지막 \code{ldexp}의 지수를 $\pm4000$으로 자르는 것은 정수 넘침을 막을 뿐 결과를 바꾸지 않는다. binary64의 지수 범위는 그보다 훨씬 좁아 $\pm4000$을 넘는 $e$는 어차피 $0$이나 $\infty$이기 때문이다. 복소수에서는 실수부와 허수부 가운데 큰 쪽의 지수로 둘을 함께 나눈다.

행렬식이 정말로 범위를 벗어나는 응용(다변량 정규 분포의 로그 밀도처럼)에는 \code{slogdet}이 있다. $\log|\det A|=\sum_i\log|u_{ii}|$와 부호(복소수면 $\prod_iu_{ii}/|u_{ii}|$에 교환의 부호를 곱한 단위 복소수)를 따로 돌려준다. 로그의 합은 넘치지 않으며, 각 로그의 상대오차가 $u$ 정도라 합의 절대오차는 많아야 $u\sum_i|\log|u_{ii}||$에 $n$번의 덧셈 오차를 더한 정도이다. 특이 행렬이면 \code{det}은 $0$, \code{slogdet}은 $(0,-\infty)$이다.
\section{From Arrays to LAPACK}
\label{sec:lu-packing}

\subsection{The Work Buffer}

stardust의 배열은 임의의 바이트 간격(stride)을 가진다. C 순서(행 우선)가 기본이고, 전치하거나 잘라 낸 뷰는 열 우선이거나 띄엄띄엄하며, 간격이 음수이거나($x[::-1]$) $0$(브로드캐스트된 축)일 수도 있다. LAPACK은 열 우선으로 연속하고 열 사이 간격이 leading dimension $\mathit{ld}$인 행렬만 받는다. 게다가 \code{getrf}는 입력을 덮어쓴다. 그래서 LU 계열의 모든 함수는 행렬을 먼저 작업 버퍼에 옮기며, 이 일을 \term{패킹}{packing}이라 한다. 패킹은 한 번의 패스로 세 가지를 한다. (i) 임의의 간격으로 읽어 열 우선으로 쓴다. (ii) \code{bool}과 \code{int64}를 \code{float64}로 바꾼다. (iii) NaN과 $\pm\infty$가 있는지 본다. LAPACK은 유한하지 않은 입력에 대해 아무것도 약속하지 않으므로(무한 루프나 쓰레기 값이 날 수 있다) stardust는 LAPACK을 부르기 전에 반드시 검사한다.

\begin{algorithm}[t]
\begin{algorithmic}[1]
\Require 원본 시작 $p$, 크기 $m\times n$, 행·열 간격 $r_s$, $c_s$ (바이트), 목적지 $D$와 $\mathit{ld}$
\Ensure $D[j\cdot\mathit{ld}+i]=x_{ij}$ (변환 포함); 모든 원소가 유한한지
\If{$r_s=$ 원소 크기} \Comment{열이 연속: 열마다 복사}
  \For{$j=0,\dots,n-1$}
    \State 열 $j$를 복사하면서 비트로 검사: $\mathit{bad}\gets\mathit{bad}\lor[(w\mathbin{\&}\mathtt{0x7ff0\ldots0})=\mathtt{0x7ff0\ldots0}]$
  \EndFor
\ElsIf{$c_s=$ 원소 크기} \Comment{행이 연속(C 순서): 타일 전치}
  \For{행 타일 $i_0=0,32,64,\dots$} \For{열 타일 $j_0=0,32,64,\dots$}
    \For{$j\in[j_0,j_0+32)$} \For{$i\in[i_0,i_0+32)$}
      \State $D[j\cdot\mathit{ld}+i]\gets x_{ij}$
    \EndFor\EndFor
  \EndFor\EndFor
  \State 써 넣은 열마다 비트 검사를 한 번 더 돈다
\Else \Comment{그 밖: 원소마다}
  \State 열마다, 행마다 읽고 검사하며 쓴다
\EndIf
\State \Return \textbf{not} $\mathit{bad}$
\end{algorithmic}
\caption{패킹 (stardust \code{pack}의 세 경로, 실수 \code{float64}의 경우)}
\label{alg:lu-pack}
\end{algorithm}

\cref{alg:lu-pack}의 경로마다 고른 이유가 있다.
\begin{itemize}
\item \textbf{비트로 하는 유한성 검사.} binary64에서 NaN과 $\pm\infty$는 지수 11비트가 모두 $1$인 수이고 그 밖에는 없다(10장). 그래서 64비트 정수로 읽어 지수 마스크를 비교하고 그 결과를 \code{or}로 모은다. 정수의 \code{or}는 결합적이어서 컴파일러가 루프를 SIMD로 벡터화한다. 부동소수점 비교와 \code{and}로 쓰면 결과는 같지만, stardust가 켠 \code{-ffp-contract=off} 아래에서 컴파일러가 이것을 순서가 정해진 축약으로 보아 벡터화하지 않았다.
\item \textbf{C 순서의 타일 전치.} 행 우선 원본을 열 순서로 그냥 읽으면 원소 하나를 읽을 때마다 행 하나($8n$바이트)를 건너뛴다. $1000\times1000$이면 원소마다 다른 캐시 줄이다. $32\times32$ 타일 하나는 읽는 쪽 32줄과 쓰는 쪽 32줄, 합쳐 $16$KB만 건드리므로 L1 캐시에 머문다(\cref{fig:lu-tile}). 원소가 $4096$개를 넘으면 타일 안을 다시 $4\times4$ 블록으로 나누어 NEON 레지스터 넷에 행 넷을 읽고 \code{trn1}·\code{trn2}로 열 넷을 만든다. 원소 루프의 약 $1.5$배이다. 원소가 $2^{17}$개를 넘는 C 순서 행렬은 열 $32$개 단위로 스레드 풀에 나누어 전치한다. 각 태스크가 자기 열만 쓰므로 결과는 직렬과 같다.
\item \textbf{어느 경로든 같은 결과.} 원소 $(i,j)$는 언제나 $D[j\cdot\mathit{ld}+i]$로 가고 변환은 원소마다 같은 함수이므로, 입력의 배치(C, F, 띄엄띄엄, 음수 간격, 브로드캐스트)가 무엇이든 작업 버퍼는 바이트까지 같다. 그래서 패킹을 거치는 \code{lu}, \code{solve}, \code{inv}, \code{det}, \code{slogdet}의 결과는 입력의 메모리 배치와 무관하게 비트까지 같다.
\end{itemize}

\begin{figure}[t]
\centering
\begin{tikzpicture}[x=0.32cm,y=0.32cm,font=\small]
\draw (0,0) rectangle (12,12); \node at (6,-1) {원본 (행 우선)};
\fill[gray!40] (4,4) rectangle (8,8);
\foreach \r in {4.5,5.5,6.5,7.5}{\draw[->] (4.2,\r) -- (7.8,\r);}
\draw (15,0) rectangle (27,12); \node at (21,-1) {작업 버퍼 (열 우선)};
\fill[gray!40] (19,4) rectangle (23,8);
\foreach \c in {19.5,20.5,21.5,22.5}{\draw[->] (\c,7.8) -- (\c,4.2);}
\draw[->,thick] (8.6,6) -- (18.4,6) node[midway,above] {타일 하나 $32\times32$};
\end{tikzpicture}
\caption{C 순서 원본의 타일 전치. 원본에서는 타일의 행을 따라(화살표) 연속으로 읽고 버퍼에서는 열을 따라 연속으로 쓴다. 한 타일이 건드리는 캐시 줄은 양쪽 합쳐 64줄이다.}
\label{fig:lu-tile}
\end{figure}

패킹은 출력 자리를 작업 버퍼로 쓰기도 한다. \code{solve}는 우변 $b$를 처음부터 결과 배열의 자리에 열 우선으로 패킹하고, 그 자리에서 \code{getrs}(또는 대입)가 $b$를 해로 덮어쓴다. 행렬 $a$만 따로 작업 버퍼가 필요하다. 그래서 \code{solve} 한 번의 메모리 이동은 ``$a$의 복사 한 번, $b$의 복사 한 번, LAPACK''이 전부이다.

\subsection{Leading Dimension and Alignment}

작업 버퍼의 leading dimension은 $\mathit{ld}=n$이다. \cref{sec:lu-blocked}에서 직접 짠 블록 LU가 $\mathit{ld}$가 2의 거듭제곱일 때 느려지는 것을 보았는데 Accelerate의 \code{dgetrf}도 같다. $n=1024$에서 $\mathit{ld}=1024$이면 초당 $149\times10^9$ 플롭, $\mathit{ld}=1032$이면 $235\times10^9$ 플롭이다(\code{figures/48/lda2.c}). stardust는 $\mathit{ld}=n$으로 패킹하므로 $1024\times1024$의 \code{solve}는 이 손해를 그대로 본다. 버퍼의 $\mathit{ld}$를 몇 원소 늘리면 피할 수 있지만 그러면 결과의 끝 비트가 바뀔 수 있어 \cref{sec:lu-testing}의 비트 관계를 함께 다시 확인해야 한다.

정렬은 더 미묘하다. Accelerate는 배열이 16바이트 경계에서 시작하는지에 따라 다른 코드 경로를 타고, 그러면 반올림이 달라진다. \code{float64}의 $3\times3$ 행렬들이 한 배열에 이어 있으면 행렬 사이 간격이 $72$바이트여서 홀수 번째 행렬은 16바이트 경계에서 8바이트 어긋난다(\cref{fig:lu-align}). 같은 행렬이 배치 안의 위치에 따라 다른 비트를 받는 것이다. stardust의 설계 문서는 $n=3,5,7,9$의 무작위 행렬 $120$개에서 \code{schur}, \code{svd}, \code{eigh}의 결과가 수십 개씩 달랐다고 적고 있다. 그래서 LAPACK이나 BLAS에 넘기는 모든 배열을 16바이트 경계에 둔다. 출력 자리가 어긋나 있으면 \code{LapackSlot}이 정렬된 임시 버퍼를 주고, 계산이 끝나면 그 내용을 출력 자리로 복사한다. 그 결과 배치의 어느 자리에 있든 같은 행렬은 같은 비트를 받는다.

\begin{figure}[t]
\centering
\begin{tikzpicture}[x=0.031cm,y=0.5cm,font=\small]
\foreach \k in {0,...,4}{
  \pgfmathsetmacro{\xa}{72*\k}
  \pgfmathtruncatemacro{\odd}{mod(\k,2)}
  \ifnum\odd=1 \fill[gray!35] (\xa,0) rectangle (\xa+72,1); \else \fill[gray!10] (\xa,0) rectangle (\xa+72,1); \fi
  \draw (\xa,0) rectangle (\xa+72,1);
  \node at (\xa+36,0.5) {행렬 \k};
  \node[below] at (\xa,0) {\pgfmathprintnumber{\xa}};
}
\foreach \b in {0,16,...,352}{\draw[gray] (\b,1) -- (\b,1.35);}
\node[anchor=west] at (0,1.75) {16바이트 경계(위의 눈금)};
\end{tikzpicture}
\caption{$3\times3$ \code{float64} 행렬 다섯 개가 이어진 배치. 행렬은 $72$바이트마다 시작하므로 홀수 번째(진한 칸)는 16바이트 경계에서 8바이트 어긋난다.}
\label{fig:lu-align}
\end{figure}

\section{The LU Family in stardust}
\label{sec:lu-stardust-family}

\subsection{Factorization: \texorpdfstring{\code{lu}}{lu}}

\code{sd.linalg.lu(x)}는 $m\times n$ 행렬(과 그 배치)을 받아 \code{LUResult(L, U, perm)}을 돌려준다. $p=\min(m,n)$이면 $L$은 $m\times p$, $U$는 $p\times n$이고 \code{x[..., perm, :] == L @ U}이다. 원소 하나의 몸체는 패킹, \code{getrf}(크기에 따라 \code{getf2}, \cref{sec:lu-batches}), 풀어 쓰기, 순열 만들기의 넷이다. \code{getrf}가 돌려주는 \code{info}$>0$(정확히 $0$인 피벗)은 무시한다. \cref{sec:lu-factorization}에서 보았듯 특이 행렬도 분해는 있고 그 흔적은 $U$의 대각에 남기 때문이다.

풀어 쓰기는 \code{getrf}가 한 자리에 섞어 둔 $L$과 $U$를 두 결과 배열로 나누는 일이다. $L$에는 대각 위에 $0$, 대각에 $1$, 대각 아래에 승수를 쓰고 $U$에는 대각과 그 위를 쓰고 아래에 $0$을 쓴다. 원소가 $1024$개 이하이면 원소마다 조건을 보는 한 루프로 하고, 그보다 크면 열마다 ``$0$ 채우기, $1$, 복사''를 \code{fill}과 \code{copy}로 한다. 작은 행렬에서는 라이브러리 호출 몇 번의 고정비가 원소 루프보다 비싸고, 큰 행렬에서는 연속 구간의 복사가 훨씬 빠르기 때문이다. 마지막으로 \cref{alg:lu-perm}으로 \code{ipiv}를 \code{perm}으로 바꾼다.

\subsection{Solving: \texorpdfstring{\code{solve}}{solve} and \texorpdfstring{\code{lu\_solve}}{lu\_solve}}

\code{solve(a, b)}는 언제나 LU로 푼다. scipy의 \code{assume\_a=}처럼 행렬의 구조를 보고 방법을 고르는 일은 하지 않는다. 대칭 양정치 행렬은 \code{cholesky\_solve}, 삼각 행렬은 \code{solve\_triangular}처럼 구조는 함수 이름으로 고르는 것이 stardust의 규칙이다. $b$가 1차원이면 벡터 하나이고 그 밖에는 $(\dots,n,k)$의 행렬 묶음이다.

\begin{algorithm}[t]
\begin{algorithmic}[1]
\Require 배치 원소 하나의 $a$ ($n\times n$), $b$ ($n\times k$)
\Ensure 결과 자리에 $x=a^{-1}b$, 또는 실패 코드
\State $a$를 작업 버퍼에 패킹; 유한하지 않으면 \Return \code{NonFinite}($a$)
\State $\mathit{info}\gets$ \code{getrf} 또는 \code{getf2}$(a)$ \Comment{$32\le n\le46$이면 \code{getf2}}
\State \textbf{if} $\mathit{info}>0$ \textbf{then return} \code{Singular}
\State $b$를 결과 자리에 패킹; 유한하지 않으면 \Return \code{NonFinite}($b$)
\If{$k\le2$}
  \State 결과의 열마다 \code{ipiv}의 교환을 적용
  \State \code{trsv}(단위 하삼각 $L$), \code{trsv}($U$)를 열마다
\Else
  \State \code{getrs}('N') \Comment{교환 + \code{trsm} 둘}
\EndIf
\end{algorithmic}
\caption{\code{solve}의 배치 원소 하나}
\label{alg:lu-solve-body}
\end{algorithm}

\cref{alg:lu-solve-body}의 순서는 에러의 순서이기도 하다. $a$가 NaN을 품고 있으면 $b$에 무엇이 있든 $a$에 대한 \code{ValueError}이고, $a$가 특이하면 $b$를 보기 전에 \code{LinAlgError}이다. 6--7행이 \code{getrs} 대신 \code{trsv}를 부르는 이유는 측정이다. Accelerate의 \code{trsm}은 우변이 하나일 때 \code{trsv}보다 크게 느리다. \cref{tab:lu-trsv}처럼 $n=1600$에서 세 배이다. \code{getrs}는 안에서 \code{trsm}을 부르므로 우변이 적을 때는 교환과 \code{trsv} 두 번을 직접 하는 편이 빠르다. 우변이 넷 이상이면 \code{trsm}이 블록으로 재사용을 얻어 앞서고, 경계는 $k\le2$에 두었다.

\begin{table}[t]
\begin{sdtable}{colspec={rrrrrrr}, row{1}={font=\small\bfseries}}
\toprule
$n$ & $50$ & $100$ & $200$ & $400$ & $800$ & $1600$ \\
\midrule
\code{trsv} & $0.42$ & $1.50$ & $5.04$ & $8.54$ & $37.4$ & $112$ \\
\code{trsm} & $0.54$ & $3.33$ & $13.0$ & $22.3$ & $100$ & $342$ \\
\bottomrule
\end{sdtable}
\caption{우변 하나의 하삼각 풀이(Accelerate, 마이크로초, \code{figures/48/bench.c}의 \code{trsv}).}
\label{tab:lu-trsv}
\end{table}

\code{lu\_solve(l, u, perm, b, *, trans='N')}는 \code{lu}의 결과를 다시 받아 푼다. scipy의 \code{lu\_factor}가 돌려주는 \code{(lu, piv)} 같은 압축 형식이 아니라 \code{lu}의 결과 그대로를 받는다는 점이 다르고, 그래서 \code{perm}이 정말 순열인지부터 확인해야 한다.

\begin{algorithm}[t]
\begin{algorithmic}[1]
\Require $l$ ($n\times n$, 대각 아래만 읽음), $u$ ($n\times n$, 대각과 위만 읽음), $\mathit{perm}$, $b$ ($n\times k$), $\mathit{trans}\in\{\mathrm N,\mathrm T,\mathrm H\}$
\State $l$의 대각 아래가 유한한지; 아니면 \Return \code{NonFinite}($l$)
\State $u$의 대각과 위가 유한한지; 아니면 \Return \code{NonFinite}($u$)
\State $u$의 대각에 정확한 $0$이 있으면 \Return \code{Singular}
\State $\mathit{seen}_{0:n}\gets0$
\For{$r=0,\dots,n-1$} \Comment{순열 검증}
  \State $v\gets\mathit{perm}_r$; \textbf{if} $v<0$ \textbf{or} $v\ge n$ \textbf{or} $\mathit{seen}_v$ \textbf{then return} \code{BadPerm}
  \State $\mathit{seen}_v\gets1$
\EndFor
\If{$\mathit{trans}=\mathrm N$}
  \State $x_{r,:}\gets b_{\mathit{perm}_r,:}$ (유한성 검사 포함) \Comment{행 모으기: 반올림 없음}
  \State $x\gets L^{-1}x$ (단위 대각); $x\gets U^{-1}x$
\Else
  \State $z\gets b$ (유한성 검사 포함)
  \State $z\gets\mathrm{op}(U)^{-1}z$; $z\gets\mathrm{op}(L)^{-1}z$ \Comment{$\mathrm{op}$은 전치 또는 켤레 전치}
  \State $x_{\mathit{perm}_r,:}\gets z_{r,:}$ \Comment{흩기}
\EndIf
\end{algorithmic}
\caption{\code{lu\_solve}의 배치 원소 하나}
\label{alg:lu-lusolve}
\end{algorithm}

\cref{alg:lu-lusolve}의 전치 경우가 맞는 까닭은 이렇다. \code{perm}의 뜻에서 $PA=LU$이고 $(PA)_r=A_{\mathit{perm}_r}$이다. 그러면 $A=P^{\mathsf T}LU$이고 $A^{\mathsf T}=U^{\mathsf T}L^{\mathsf T}P$이다. $A^{\mathsf T}x=b$는 $U^{\mathsf T}w=b$, $L^{\mathsf T}z=w$, $Px=z$로 풀리고 마지막 식은 $x_{\mathit{perm}_r}=z_r$, 곧 흩기이다. \code{'N'}이면 반대로 $Pb$를 먼저 만드는데 이것은 $b$의 행을 \code{perm}의 순서로 모으는 것이어서 산술이 없다. 그 뒤의 두 삼각 풀이는 \code{solve}가 교환 뒤에 하는 것과 같은 호출이다. 그래서 $l$, $u$가 \code{lu}가 돌려준 열 우선 배열 그대로이면 \code{lu\_solve(*lu(a), b)}는 \code{solve(a, b)}와 비트까지 같다.

$l$과 $u$는 LAPACK이 읽기만 하므로 복사하지 않고 그 자리에서 읽을 수 있다. 원소 타입이 계산 타입과 같고, 시작이 16바이트 경계이고, 열 우선이면서 열 간격이 16바이트의 배수이면 그대로 \code{trsv}·\code{trtrs}에 넘기고 읽을 삼각형만 유한성을 검사한다. C 순서(행 우선)인 삼각 행렬도 그 자리에서 읽는다. 행 우선의 하삼각은 열 우선으로 본 상삼각의 전치이므로, 우변이 둘 이하이면 \code{trsv}에 \code{CblasRowMajor}를 주고 셋 이상이면 \code{trtrs}에 반대쪽 삼각과 반대 연산(\code{'N'}$\leftrightarrow$\code{'T'})을 준다. 다만 이 길은 반올림 순서가 달라 열 우선 인수와 끝 비트가 다를 수 있다. 복소수의 켤레 전치(\code{'H'})에 우변이 셋 이상인 경우는 대응하는 연산이 없어서 패킹한다.

\subsection{Inverse and Determinants}

\code{inv(x)}는 \code{getrf} 뒤에 단위행렬을 우변으로 두고 \cref{alg:lu-solve-body}의 풀이를 그대로 한다. \code{getri}를 쓰지 않는 이유는 둘이다. 하나는 \cref{tab:lu-inverse}에서 보았듯 $n\le1024$에서 이 길이 빠르거나 같다는 것이고, 다른 하나는 \code{inv(a)}가 \code{solve(a, I)}와 비트까지 같아진다는 것이다. 두 함수가 같은 행렬에 대해 반올림 수준에서 다른 답을 주면 사용자가 의아해한다. 대신 $n=2048$처럼 큰 행렬에서는 \code{getri}보다 느리다. 설계 문서는 크기에 따라 길을 바꾸는 안을 남겨 두었다.

\code{det}과 \code{slogdet}는 \code{getrf} 뒤에 \cref{alg:lu-det}과 로그 합을 한다. 특이(\code{info}$>0$)도 에러가 아니다. 빈 행렬($0\times0$)의 행렬식은 빈 곱의 관례대로 $1$, \code{slogdet}은 $(1,0)$이다. 배치가 없으면 결과는 Python 스칼라이고 배치가 있으면(크기가 1이어도) 배열이다.

\subsection{Errors}

LU 계열이 내는 에러는 둘이다. 입력이 NaN이나 무한대를 품으면 \code{ValueError}이고 특이 행렬로 풀거나 역행렬을 구하면 \code{LinAlgError}이다. 메시지는 함수 이름으로 시작한다.
\begin{itemize}
\item \code{solve(): 'a' contains NaN or infinity}
\item \code{inv(): singular matrix (at batch index (1, 2))}
\end{itemize}
배치의 여러 원소가 실패하면 가장 작은 배치 인덱스를 보고한다. 한 원소 안에서는 인자의 순서대로 본다. 첫 행렬의 유한성, 그 행렬의 특이성, 다음 인자들의 유한성이다(\cref{alg:lu-solve-body,alg:lu-lusolve}). ``특이''는 정확히 $0$인 피벗만을 뜻한다. 피벗이 $10^{-300}$이어도 분해는 성공하고 답은 거대하거나 무한대가 되며, 조건수를 보는 것은 사용자의 몫이다(\code{cond}). 이 규칙은 numpy와 같다.

\section{Batches, Threads, and Hardware}
\label{sec:lu-batches}

\subsection{The Batch Driver}

모든 행렬 함수는 앞쪽 축을 \term{배치}{batch}로 본다. \code{solve(a, b)}에서 $a$가 $(2,1,n,n)$이고 $b$가 $(3,n,k)$이면 배치 축 $(2,1)$과 $(3,)$을 오른쪽 정렬로 브로드캐스트해 $(2,3)$개의 문제를 푼다. 원소마다 각 피연산자의 바이트 오프셋을 계산하고(브로드캐스트되는 축의 간격은 $0$) 같은 몸체를 부른다. 배치 원소는 서로 독립이므로 스레드에 나누어 돌릴 수 있다.

\begin{algorithm}[t]
\begin{algorithmic}[1]
\Require 배치 크기 $N$, 원소 하나의 추정 시간 $t$ (ns), 몸체 $\mathit{body}(i)$ (실패 코드 또는 $0$), 태스크 상한 $c$
\If{$N\ge2$ \textbf{and} $Nt\ge10^4$ \textbf{and} 풀이 쓸 수 있음}
  \State $q\gets\max(1,\lfloor5000/t\rfloor)$ \Comment{태스크 하나가 약 $5\,\mu$s}
  \State \textbf{if} $c>0$ \textbf{and} $\lceil N/q\rceil>c$ \textbf{then} $q\gets\lceil N/c\rceil$
  \State $\mathit{first}\gets\infty$ (원자적)
  \ForAll{태스크 $\tau=0,\dots,\lceil N/q\rceil-1$ \textbf{in parallel}}
    \State 자기 작업 메모리를 만든다
    \For{$i=\tau q,\dots,\min(N,(\tau+1)q)-1$}
      \State \textbf{if} $i>\mathit{first}$ \textbf{then break} \Comment{더 작은 실패가 이미 있다}
      \If{$\mathit{body}(i)\ne0$}
        \State 기록 $(i,\text{코드})$; $\mathit{first}\gets\min(\mathit{first},i)$ (비교 후 교환); \textbf{break}
      \EndIf
    \EndFor
  \EndFor
  \State \Return 기록들 가운데 인덱스가 가장 작은 것
\Else
  \State GIL을 풀고(오래 걸리면) $i=0,1,\dots$ 차례로 $\mathit{body}(i)$; 첫 실패에서 멈추고 \Return
\EndIf
\end{algorithmic}
\caption{배치 드라이버 (\code{run\_batch})}
\label{alg:lu-batch}
\end{algorithm}

\cref{alg:lu-batch}에서 원소 하나의 추정 시간은 $t=a+c\,w$의 꼴이다. $w$는 행렬의 원소 수이고 LU 계열은 $a=70$ns, $c=3$ns로 $n=2..256$의 실측에 맞춘 상수이다. $n$이 크면 시간이 $n^3$으로 자라 이 식은 과소평가하지만, 그때는 원소 둘만 있어도 이미 나눌 가치가 있으므로 판단은 틀리지 않는다. 배치 전체의 추정 시간이 $10\,\mu$s 미만이면 스레드를 깨우는 비용이 이득보다 크므로 나누지 않는다.

실패 보고는 직렬과 같아야 한다. 직렬로 돌면 처음 만난 실패가 가장 작은 인덱스이다. 나누어 돌릴 때 처음 발견한 실패에서 모두 멈추면 어느 인덱스가 보고될지가 스레드의 타이밍에 달린다. \cref{alg:lu-batch}은 이렇게 한다. 각 태스크는 자기 구간을 앞에서부터 돌다 첫 실패에서 멈추고 그 인덱스를 기록한다. 지금까지 알려진 가장 작은 실패보다 뒤의 원소는 건너뛴다. 끝나면 기록들 가운데 가장 작은 것을 고른다. 가장 작은 실패 인덱스 $i^\ast$를 품은 태스크는 $i^\ast$에 이르기 전에 멈추지 않는다. 그 앞의 원소들은 실패하지 않으므로 기록되지 않고, $\mathit{first}$는 언제나 어떤 실제 실패의 인덱스여서 $i^\ast$ 이상이므로 건너뛰기 조건도 걸리지 않는다. 그러므로 $i^\ast$는 반드시 기록되고, 보고되는 인덱스는 스레드 수와 상관없이 직렬과 같다.

결과의 비트도 스레드 수와 상관없다. 원소 하나는 어디서 돌든 같은 몸체를 같은 순서로 돌고, 작업 메모리는 태스크마다 따로이며, LAPACK에 넘기는 배열은 모두 16바이트 경계에 있기 때문이다. 작업 메모리에도 두 가지 세부가 있다. 하나는 새 메모리를 $0$으로 채우지 않는 것이다. \code{std::vector}의 \code{resize}처럼 채우면 $1000\times1000$에서 그 비용이 패킹만큼 들었고, 모든 사용처가 읽기 전에 쓰므로 채울 필요가 없다. 다른 하나는 $4$KB보다 작은 버퍼를 $128$바이트 경계(이 처리기의 캐시 줄)에 두고 크기도 $128$의 배수로 올리는 것이다. 그러지 않으면 메모리 할당기가 여러 스레드의 작은 버퍼(\code{lu\_solve}의 $3$바이트짜리 \code{seen})를 같은 캐시 줄에 내주는 일이 생겼고, 원소마다 쓰는 그 줄이 코어 사이를 오가며 시간이 프로세스마다 30\%씩 달라졌다.

\subsection{A Shared Matrix Unit}

Apple의 처리기에는 행렬 곱을 위한 장치(AMX)가 코어 무리마다 하나씩 있고 Accelerate의 일부 경로가 이것을 쓴다. 이 장치는 코어들이 함께 쓰므로 여러 스레드가 동시에 그 경로를 타면 줄을 선다. \cref{fig:lu-amx}는 크기가 같은 행렬을 수천 개 분해하는 시간을 행렬 하나당으로 나타낸 것이다. 한 스레드에서는 \code{getrf}와 \code{getf2}가 비슷하다. 열 스레드에서 \code{getf2}는 $n$을 따라 고르게 늘지만 \code{getrf}는 $n=32$에서 갑자기 $1.5\,\mu$s로 뛰고 $n=64$까지 그 자리에 머문다. 스레드 수를 늘려도 빨라지지 않는 이 평평한 구간이 공유 장치에서 줄을 서는 모습이다. $n=36$에서 \code{getf2}는 \code{getrf}의 $4.2$배 빠르다.

\begin{figure}[t]
\centering
\begin{tikzpicture}
\begin{axis}[width=0.86\textwidth,height=0.5\textwidth,
  xlabel={$n$},ylabel={행렬 하나당 시간 ($\mu$s)},grid=major,grid style={gray!25},
  legend pos=north west,legend style={font=\small,draw=none,fill=none},
  xmin=6,xmax=74,ymin=0,ymax=2.5]
\addplot[black,mark=*,mark size=1.2pt] table[x=n,y=getrf10]{figures/48/batch.dat};
\addlegendentry{\code{getrf}, 10스레드}
\addplot[black,dashed,mark=square*,mark size=1.2pt] table[x=n,y=getf210]{figures/48/batch.dat};
\addlegendentry{\code{getf2}, 10스레드}
\draw[gray!60] (axis cs:32,0) -- (axis cs:32,2.5);
\draw[gray!60] (axis cs:46,0) -- (axis cs:46,2.5);
\node[font=\scriptsize,anchor=south] at (axis cs:39,2.05) {stardust의 \code{getf2} 구간};
\end{axis}
\end{tikzpicture}
\caption{$n\times n$ 행렬 수천 개를 열 스레드로 분해할 때 행렬 하나당 시간(Apple M4 Pro, Accelerate). 한 스레드의 값과 측정 코드는 \code{figures/48/batch.dat}과 \code{bench.c}에 있다.}
\label{fig:lu-amx}
\end{figure}

stardust는 이 측정을 두 규칙으로 바꿨다. 첫째, 행과 열이 모두 $32$에서 $46$ 사이인 행렬은 \code{getrf} 대신 블록 없는 \code{getf2}로 분해한다. 피벗 규칙과 특이의 판정은 같고 결과는 반올림 수준에서 같다. 둘째, 우변이 $8$개 이상인 삼각 풀이(\code{trsm})도 $n=24$--$40$에서 같은 장치를 타므로 이 경우 배치를 나누는 태스크의 수를 $3$이나 $5$로 제한한다(\cref{alg:lu-batch}의 $c$). 경계값은 stardust를 개발한 기계에서 잰 것이다. 우리 측정에서는 열 스레드의 \code{getf2}가 $n=64$까지 앞서므로 경계는 처리기마다 다르다. 이런 규칙은 하드웨어를 바꿀 때마다 다시 재야 한다는 점에서 알고리즘의 일부가 아니라 구현의 일부이다.

\section{Testing and Accuracy}
\label{sec:lu-testing}

후진 오차 해석은 시험의 기준이 된다. stardust는 풀이의 \term{잔차 계약}{residual contract}을
\[
\|Ax-b\|_\infty\le c\,n\,\epsilon\,(\|A\|_\infty\|x\|_\infty+\|b\|_\infty),\qquad c=10
\]
으로 둔다. \cref{thm:lu-rigal-gaches}의 후진 오차를 $\infty$-노름으로 쓴 것에 $c\,n\,\epsilon$의 상한을 준 꼴이다. \eqref{eq:lu-normwise}의 이론 상한 $n^2\rho_n\gamma_{3n}$은 너무 비관적이어서 시험 기준으로 쓸 수 없고, 상수 $c=10$은 측정에서 정했다. \cref{tab:lu-residual}은 무작위 정규 행렬에서 잰 이 비율이다. 가장 나쁜 값도 계약의 $1/90$ 아래이다.

\begin{table}[t]
\begin{sdtable}{colspec={rrrr}, row{1}={font=\small\bfseries}}
\toprule
$n$ & 행렬 수 & 중앙값 & 최댓값 \\
\midrule
$10$ & $2000$ & $0.019$ & $0.112$ \\
$100$ & $300$ & $0.0055$ & $0.0092$ \\
$1000$ & $20$ & $0.0036$ & $0.0049$ \\
\bottomrule
\end{sdtable}
\caption{$\|Ax-b\|_\infty/(n\epsilon(\|A\|_\infty\|x\|_\infty+\|b\|_\infty))$, 표준 정규 원소의 $A$와 $b$, 부분 피벗 LU(Accelerate)로 푼 해(\code{figures/48/residual.py}). 계약은 이 값이 $10$ 이하일 것을 요구한다.}
\label{tab:lu-residual}
\end{table}

이 밖에 stardust의 시험은 다음을 확인한다.
\begin{itemize}
\item \textbf{조건수가 큰 행렬}: $\kappa_2=1$에서 $10^{12}$까지의 정사각 행렬에 대해 해와 행렬식의 참값을 구간 산술 라이브러리 Arb로 미리 계산해 두고, 잔차 계약과 함께 확인한다.
\item \textbf{비트 관계}: \code{inv(a) == solve(a, I)}, \code{lu\_solve(*lu(a), b) == solve(a, b)}(열 우선 인수), \code{matrix\_power(x, -1) == inv(x)}가 바이트까지 같은지.
\item \textbf{배치와 레이아웃}: 같은 행렬을 C, F, 띄엄띄엄, 음수 간격으로 넣었을 때, 배치의 홀수 자리($n=3,5,7,40$)에 두었을 때, 스레드를 $1$, $3$, $10$개로 바꿨을 때 결과가 같은 바이트인지.
\item \textbf{경계 크기}: \code{getf2} 구간의 양 끝 $n=31,32,46,47$에서 답, 비트 재현성, 특이의 판정.
\item \textbf{규칙}: 빈 행렬, 배치 크기 $0$, NaN 입력의 \code{ValueError}, 특이 행렬의 \code{LinAlgError}와 가장 작은 배치 인덱스, 에러의 순서, 넘치는 부분 곱이 있는 \code{det}(\code{diag(1e200, 1e200, 1e-300)}이 $10^{100}$), 복소 \code{lu}의 $|\ell_{ij}|\le\sqrt2$.
\end{itemize}

numpy와의 비교는 허용오차로 한다. Accelerate의 \code{gesv}는 \code{getrf}와 \code{getrs}를 따로 부를 때와 내부 계산 순서가 달라 $n\ge4$에서 끝 비트가 다르고, numpy의 \code{solve}는 \code{gesv}를 부르기 때문이다. 같은 기계와 같은 스레드 수에서 실행할 때마다 같은 비트를 내는 것이 stardust의 약속이고, 다른 라이브러리와 비트가 같은 것은 약속이 아니다.

\section{APIs in This Chapter}
\label{sec:lu-apis}

\begin{table}[h]
\begin{sdtable}{colspec={lXl}, row{1}={font=\small\bfseries}}
\toprule
API & 하는 일 & 설명한 곳 \\
\midrule
\code{lu(x)} & $x[\mathit{perm}]=LU$, 직사각·특이 허용 & \cref{sec:lu-pivoting,sec:lu-blocked,sec:lu-stardust-family}, \cref{alg:lu-blocked,alg:lu-recursive,alg:lu-perm} \\
\code{solve(a, b)} & $ax=b$, 배치·브로드캐스트 & \cref{sec:lu-using,sec:lu-stardust-family}, \cref{alg:lu-solve,alg:lu-solve-body} \\
\code{lu\_solve(l, u, perm, b, *, trans)} & \code{lu}의 결과로 풀기, 전치 & \cref{sec:lu-stardust-family}, \cref{alg:lu-lusolve} \\
\code{inv(x)} & 역행렬, \code{solve(x, I)}와 같은 비트 & \cref{sec:lu-using,sec:lu-stardust-family}, \cref{tab:lu-inverse} \\
\code{det(x)} & 넘침을 피하는 행렬식 & \cref{sec:lu-using}, \cref{alg:lu-det} \\
\code{slogdet(x)} & 부호와 $\log|\det|$ & \cref{sec:lu-using} \\
\bottomrule
\end{sdtable}
\caption{이 장이 구현까지 설명한 \code{sd.linalg}의 API. 패킹과 배치 드라이버(\cref{sec:lu-packing,sec:lu-batches})는 이 장의 모든 API가 함께 쓴다.}
\label{tab:lu-apis}
\end{table}

\section{Concluding Remarks}

가우스 소거를 행렬의 분해로 보는 순간 알고리즘은 세 층으로 갈린다. 산술의 층에서는 승수를 제자리에 적는 것만으로 $L$이 생기고, 부분 피벗과 성장 인자가 안정성을 정한다. 이론 상한 $2^{n-1}$과 실제 성장 $n^{2/3}$ 사이의 간극은 아직 완전히 설명되지 않았다. 메모리의 층에서는 같은 산술을 \code{gemm}으로 바꾸어 데이터 하나를 여러 번 쓰게 하는 것이 속도를 열 배 넘게 바꾸며, 블록과 재귀가 그 방법이다. 기계의 층에서는 leading dimension의 모양, 16바이트 정렬, 코어들이 함께 쓰는 행렬 장치처럼 알고리즘 책에 나오지 않는 것들이 속도와 비트 재현성을 좌우한다. stardust의 LU 계열은 산술을 LAPACK에 맡기는 대신 나머지 두 층에서 패킹, 정렬, 배치 드라이버, 크기별 경로라는 형태로 일을 한다. 다음 장에서는 삼각 행렬과 띠 행렬처럼 구조가 있는 연립방정식을 다룬다.

LU와 그 오차 해석의 표준 참고서는 Higham의 책 9장과 Golub과 Van Loan의 책 3장이다\cite{higham2002accuracy,golub2013matrix}. 가우스 소거의 역사는 Grcar가 자세히 정리했다\cite{grcar2011}.

\chapterreferences
